%% =====================================================================
%% DRAFT / extended version of Section 5.  NOT the file \input by
%% ETAPS2026.tex; the shipped version is sec/section5_OTF.tex.
%%
%% This file is the same section with every proof inline and the
%% punctuality figure kept in place (~6 pages).  The shipped version
%% (~4.6 pages) states Proposition 1 and Lemmas 1-4 without proof and
%% points at the appendix; APPENDIX A below is exactly that appendix,
%% ready to \input after the bibliography.
%%
%% Design notes are at the end of the file.
%% =====================================================================

\section{The Decision Procedure}\label{sec:correct}

Section~\ref{sec:tableau} defines success by a condition on a marking of a graph that no
implementation can afford to build. We make that graph finite and identify the largest successful
marking with a greatest fixpoint (\S\ref{ss:gfp}), give a procedure deciding membership of the root
in it while building only what the search demands (\S\ref{ss:otf}), and prove that procedure sound
for TCTL and complete for the non-punctual fragment (\S\ref{ss:sound}). Throughout $\varphi$ is the
input, $n=|\FL(\varphi)|$, and $M$ the largest constant in its intervals. 

\subsection{The tableau graph and its greatest marking}\label{ss:gfp}

The calculus leaves the order of rule application open; the procedure fixes one. Let
$\FL^{\bullet}(\varphi)$ extend $\FL(\varphi)$ with the pursuit markers $\mu^\bullet,\mu^{\ast}$ and
with $\EX\chi,\AX\chi$ for $\chi$ in it, so that it holds every goal the rules produce and has size
$O(n)$, and fix a linear order on it. Then $\Expand{N}$ is empty if $N$ is closed, and otherwise
fires on the least goal of $\Gamma$ it fits the first applicable rule in the order \ruleref{Conj},
\ruleref{Disj}, \ruleref{AX}, \ruleref{Roll}, \ruleref{Cov}, \ruleref{Rsplit},
\ruleref{Runf_before}, \ruleref{Runf_in}, \ruleref{Uunf}, \ruleref{Uunfroll}; when none applies, $N$
is saturated and $\Expand{N}$ returns the \ruleref{Delay} and \ruleref{EX} children together, both
being conjunctive. Two conventions make this a function of the node: the clock a \ruleref{launch}
allocates is determined by $\mu$ rather than drawn fresh, sound because \ruleref{Cov} blocks a second
launch while a pursuit of $\mu$ is live, so $\clockSet=\{x_\mu\mid\mu$ timed in $\FL(\varphi)\}$ has
size $O(n)$; and acceptance reads clock resets and stop branches off the graph, both properties of an
edge, so each node records the clocks its rule reset and the modality it discharged.

\begin{definition}[Tableau graph]\label{def:graph}
$G(\varphi)$ is the least set of nodes containing $N_0=\tup{\{\varphi\}}{\emptyset}{\{\mathbf
0\}}{\emptyset}{\emptyset}$ and closed under $\Expand{}$, two nodes being identified when they agree
componentwise and their zones agree modulo $\cong_M$, with an edge $N\to N'$ for each
$N'\in\Expand{N}$.
\end{definition}

\begin{lemma}[Finiteness]\label{lem:finiteness}
$|G(\varphi)|\le 2^{O(n)}\cdot M^{O(n^{2})}$, and every node has $O(n)$ children.
\end{lemma}
\begin{proof}
$\Gamma$, $\Ctx$, the reset set and the discharge marker range over subsets of
$\FL^{\bullet}(\varphi)$ or of $\clockSet$; $\obligations$ is fixed by its domain; and $\Deadline$ gives each
$x_\mu$ one of the bounds $\sup\kI$, $\sup\kI-\inf\kI$ with a strictness flag. Each contributes
$2^{O(n)}$. A canonical $\cong_M$-zone over $K\le n+1$ clocks is a $(K{+}1)^{2}$ DBM with entries
$\infty$ or in $\{-M,\dots,M\}\times\{<,\le\}$, giving $(4M{+}3)^{(n+2)^{2}}=M^{O(n^{2})}$ zones.
Every rule has at most three children except \ruleref{EX}, with one per $\EX\psi\in\Gamma$.
\end{proof}

\noindent
Identification up to $\cong_M$ dispenses with a termination argument for expansion, and so lets the
calculus afford rules that re-impose their own premise: a chain of applications that would regress
forever revisits a node and closes a cycle, which crosses no \ruleref{Delay} step and is discarded by
\ref{M1}.

A marking contains no closed node, at least one child of each of its \ORn-nodes and all children of
each of its \ANDn-nodes, and is successful when every member satisfies \ref{M1} and \ref{M2}. Both
clauses quantify over paths \emph{inside} the marking, so a successful marking is a set that survives
the deletions it prescribes. For $A\subseteq G(\varphi)$ let $F(A)$ be $A$ minus:
\begin{enumerate}[start=1,label={D\arabic*)}, itemindent=1.6em, ref=(D\arabic*), itemsep=0.02em]
\item\label{D1} the closed nodes;
\item\label{D2} the \ORn-nodes with no child in $A$, and the \ANDn-nodes with a child outside $A$;
\item\label{D3} the nodes with no $A$-path to a node on a time-progressing cycle of $A$;
\item\label{D4} for each until $\mu$, the nodes pursuing $\mu$ from which no $A$-path carrying the
      pursuit throughout reaches a stop for $\mu$ taken where the zone meets $x_\mu\in\kI$.
\end{enumerate}

\begin{proposition}\label{prop:gfp}
$F$ is monotone, and $A$ is a successful marking iff $A\subseteq F(A)$. Hence the successful markings
have a greatest element $\mathcal M^{*}=\nu F$, and the tableau succeeds on $\varphi$ iff
$N_0\in\mathcal M^{*}$.
\end{proposition}
\begin{proof}
\ref{D1} ignores $A$; a node kept by \ref{D2} has its witnessing child, or all its children, in $A$
and hence in any superset; \ref{D3} and \ref{D4} ask for a path in $A$, which survives enlargement.
So $F$ is monotone on the finite lattice $2^{G(\varphi)}$, and $A\subseteq F(A)$ unfolds to the
clauses defining a successful marking. Knaster--Tarski gives a greatest post-fixpoint, equal to $\nu
F$ and to the union of all post-fixpoints, and the tableau succeeds iff one of them contains $N_0$.
\end{proof}

\subsection{The on-the-fly procedure}\label{ss:otf}

Algorithm~\ref{alg:otf} decides $N_0\in\mathcal M^{*}$ without ever holding $G(\varphi)$.
$\Expand{}$ applies the rules the first time the search reaches a node and returns the recorded
children of a $\cong_M$-equal node otherwise, so the graph exists only as far as the search has
demanded it. Tarjan's indices, used as a black box, ride along; a back edge yields a lowlink update
and nothing else, a node still on the stack having no verdict and none being guessed for it. The
third line of the loop is where the search stops at the first witness: a child settles its parent
when the parent is an \ORn-node and the child holds, or an \ANDn-node and it fails, which is what
$\val[w]=[\,v$ is $\ORn\,]$ says, and the parent retires with its remaining children unbuilt; the guard $\mathrm{top}(\Stack)=v$ makes this well defined, an earlier child
possibly having left members of $v$'s component above $v$ on the stack. A component is decided the
moment it closes, against the verdicts of the children it leaves: writing
$\Out_C=\{w\notin C\mid w\in\ch[v]$ for some $v\in C$, $\val[w]=\top\}$ for those already accepted,
$\nu F_C$ is the largest $A\subseteq C$ with $A\subseteq F(A\cup\Out_C)$, computed by deletion from
$A=C$, each round applying \ref{D1} and \ref{D2} pointwise and then the backward reachabilities
\ref{D3}, \ref{D4} over the survivors. A deletion can disconnect $A$, so the decomposition is
recomputed each round.

\begin{algorithm}[!tb]
\footnotesize
\DontPrintSemicolon
\SetAlgoLined
\SetKwProg{Fn}{procedure}{}{end}

\Fn{\Decide{$\varphi$}}{
  $\val[v]\gets{?}$ and $\idx[v]\gets\bot$ for all $v$; \ $\Stack\gets\varepsilon$; \ $\cnt\gets0$\;
  \Explore{$N_0$} \ \textbf{where} $N_0=\tup{\{\varphi\}}{\emptyset}{\{\mathbf
    0\}}{\emptyset}{\emptyset}$; \quad \Return{\SAT{} \textnormal{iff} $\val[N_0]=\top$}\;
}
\BlankLine
\Fn{\Explore{$v$}}{
  $\idx[v]\gets\low[v]\gets\cnt{+}{+}$; \ $\mathrm{push}(\Stack,v)$; \
  $\ch[v]\gets\Expand{v}$ \tcp*{memoised on $\cong_M$}
  \ForEach{$w\in\ch[v]$}{
    \lIf{$\idx[w]=\bot$}{\Explore{$w$}; \ $\low[v]\gets\min(\low[v],\low[w])$}
    \lElseIf{$w\in\Stack$}{$\low[v]\gets\min(\low[v],\idx[w])$}
    \lIf{$\val[w]=[\,v\text{ is }\ORn\,]$ \textnormal{\textbf{and}} $\mathrm{top}(\Stack)=v$}
        {$\val[v]\gets\val[w]$; \ $\mathrm{pop}(\Stack)$; \Return}
  }
  \lIf{$\low[v]=\idx[v]$}{$C\gets$ pop $\Stack$ down to $v$; \
    $\val[u]\gets[\,u\in\nu F_C\,]$ for $u\in C$}
}
\caption{On-the-fly satisfiability.}\label{alg:otf}
\end{algorithm}

\begin{lemma}\label{lem:search}
\emph{(i)} If $\val[w]$ is final and settles $v$, then $v\in\mathcal M^{*}$ iff $w\in\mathcal M^{*}$,
and leaving the remaining children of $v$ unexpanded changes $\mathcal M^{*}$ at no node.
\emph{(ii)} A strongly connected $T$ with an internal edge carries a time-progressing cycle iff some
node of $T$ is a \ruleref{Delay}-successor and every clock is reset somewhere in $T$ or unbounded
everywhere in $T$.
\emph{(iii)} If $C$ is a maximal strongly connected set with $\val[c]=[\,c\in\mathcal M^{*}\,]$ for
every child $c\notin C$ of a node of $C$, then $\mathcal M^{*}\cap C=\nu F_C$.
\end{lemma}
\begin{proof}
(i) The first claim is \ref{D2}. The explored subgraph is closed under $\Expand{}$ except at retired
nodes, so a path of $\mathcal M^{*}$ using an unexpanded child of $v$ passes through $v$; if
$v\notin\mathcal M^{*}$ there is no such path, and if $v\in\mathcal M^{*}$ then $v$ satisfies
\ref{M1}, and \ref{M2} for each pursuit active at it, so a node reaching $v$ inherits both witnesses
by prefixing the path to $v$, which carries a pursuit exactly when it is active at $v$.

(ii) $(\Leftarrow)$ Concatenating paths in $T$ gives a cycle through the \ruleref{Delay}-successor,
which therefore elapses positive time, and through one reset node of every clock $T$ resets; the rest
are unbounded throughout. $(\Rightarrow)$ A time-progressing cycle in $T$ elapses positive time,
hence crosses a \ruleref{Delay} step; and a clock reset nowhere in $T$ is forced unbounded on that
cycle, whence unbounded on all of $T$, since a path inside $T$ never resets it and its value, already
above $M$, only grows.

(iii) $C$ being maximal, a path of $\mathcal M^{*}$ from $N\in C$ either stays in $C$, and a witness
for \ref{M1} or \ref{M2} then lies inside $C$; or it leaves $C$ once and for all at some
$c\in\mathcal M^{*}\setminus C$, which by hypothesis lies in $\Out_C$ and itself satisfies \ref{M1},
and \ref{M2} for every pursuit active at it, so prefixing the path from $N$ to $c$ turns either
witness into one for $N$. Hence $\mathcal M^{*}\cap C\subseteq F(\mathcal M^{*}\cap C\cup\Out_C)$.
Conversely any $A\subseteq C$ with $A\subseteq F(A\cup\Out_C)$ extends by the nodes of
$\mathcal M^{*}$ reachable from $\Out_C$ to a post-fixpoint of $F$, which Proposition~\ref{prop:gfp}
places inside $\mathcal M^{*}$. The greatest elements coincide.
\end{proof}

\begin{theorem}[Adequacy and complexity]\label{thm:alg}
Algorithm~\ref{alg:otf} terminates and retires every node it builds with
$\val[N]=[\,N\in\mathcal M^{*}\,]$; it returns \SAT{} iff $N_0\in\mathcal M^{*}$, and runs in time
$2^{\mathrm{poly}(n)}$.
\end{theorem}
\begin{proof}
By Lemma~\ref{lem:finiteness} with memoisation the traversal is a depth-first search of a finite
graph, so it terminates. Components are popped in reverse topological order, and by the Tarjan
invariant every child of $C$ outside $C$ has retired when $C$ is popped, one still on the stack
dragging the lowlink of $C$'s root below its index; so $\Out_C$ is final and every pop meets
Lemma~\ref{lem:search}(iii). The deletion loop decreases on a finite set and each step instantiates a
clause of $F$, exact by Lemma~\ref{lem:search}(ii), so it reaches $\nu F_C=\mathcal M^{*}\cap C$, and
by Lemma~\ref{lem:search}(i) pruning preserves every verdict.

For the bound let $m$ be that of Lemma~\ref{lem:finiteness}. As $M$ occupies at most $n$ input
symbols, $\log m=O(n)+O(n^{2})\log M=\mathrm{poly}(n)$. The explored graph has at most $m$ nodes and
$O(nm)$ edges, and memoisation on $\cong_M$-classes expands each once, avoiding the doubly
exponential blow-up of an ancestor-only unfolding. A component takes at most $|C|$ rounds,
each deleting a node or stopping, of one decomposition and $O(n)$ backward reachabilities linear in
the subgraph, and the components partition the nodes, so the run costs $2^{\mathrm{poly}(n)}$.
\end{proof}

\subsection{Soundness and Completeness}\label{ss:sound}

\begin{lemma}\label{lem:zones}
Let $N_0,\dots,N_k$ be a path of $G(\varphi)$ with non-empty zones. \emph{(i)} For every $\nu_k\in
Z_k$ there are $\nu_i\in Z_i$ and delays realising every step, positive at every \ruleref{Delay}
step. \emph{(ii)} If a delay crosses no integer hyperplane $x=b$ with $b\le M$, every subformula of
$\varphi$ has constant truth on the open segment it sweeps. \emph{(iii)} A cycle of $G(\varphi)$ has
a non-Zeno unfolding iff it is time-progressing, and then one whose iterations elapse a constant
positive amount.
\end{lemma}
\begin{proof}
(i) Every zone operation a rule performs (intersection with a single-clock guard, reset,
time-successor, $\cong_M$-canonicalisation) back-propagates, non-empty precisely because the
successor valuation lies in its image~\cite{Bouyer05TimedAutomata}; thread backwards from $\nu_k$,
noting that $\Dopen Z$ and $\Dbd Z$ lie strictly ahead of $Z$ in time. (ii) No clock crosses an
integer at or below $M$, so no integer part changes and no fractional ordering wraps, while clocks
above $M$ are identified by extrapolation; the segment lies in one region, on which TCTL truth is
constant~\cite{AlurCourcoubetisDill1990,Bouyer05TimedAutomata}, and its endpoints are nodes of their
own by the integer split of \ruleref{Delay}. (iii) A cycle elapsing no time takes zero time over
infinitely many traversals; otherwise some clock is neither reset nor extrapolated above $M$, so it
keeps a bound $c\le M$, grows by exactly the delay elapsed, and caps the total at $c$. Conversely (i)
yields a traversal elapsing some $\delta>0$, and the cycle returning to the same node, no rule
distinguishes two visits and no budget depletes, so $\delta$ recurs at every iteration.
\end{proof}

\noindent
\textbf{From a marking to an automaton.} Call a node of $\mathcal M^{*}$ a \emph{state} if it is
saturated, so that its children are the successor steps; the remaining, \emph{expansion} edges
neither elapse time nor change location, and reset a clock only through a \ruleref{launch} or a
\ruleref{Ustoproll}. At most $n+1$ obligations are pending at $N$, the divergence requirement of
\ref{M1} and the until pursuits active there, and \ref{M1}, \ref{M2} give a finite path $\pi_{N,o}$
of $\mathcal M^{*}$ discharging each. A \emph{schedule} selects at every \ORn-node the child
$\pi_{N,o}$ prescribes and moves round-robin to the next pending obligation once $o$ is discharged,
so its memory, the pair (obligation, position on $\pi_{N,o}$), is finite. The witness
$\mathcal A_\varphi$ has as locations the pairs (state, schedule memory), labelled by
$\Gamma_{\mathsf{prop}}$ and carrying $\Inv(N)$, and as edges the scheduled successor steps, guarded
by the zone of the next state reached through expansion edges and resetting the clocks that segment
resets.

\begin{lemma}[Recurrence]\label{lem:recurrence}
Let $\mu=\mathrm{A}(\varphi\,\mathbf{U}_\kI\,\psi)$ with $a=\inf\kI$, $b=\sup\kI<\infty$, and let
$\rho$ be a run of $\mathcal A_\varphi$ through a node carrying $\mu^{\ast}$, launched at $t_0$. Then
\ruleref{Ustoproll} fires on $\rho$ at times $t_0<w_1<w_2<\cdots$ with $w_1-t_0\in\kI$ and
$w_{j+1}-w_j\le b-a$, and $\rho\models\varphi\,\mathbf{U}_\kI\,\psi$ from every position at or after
$t_0$.
\end{lemma}
\begin{proof}
No rule drops $\mu^{\ast}$: \ruleref{Ustoproll} keeps it and \ruleref{Ugo} moves it to $\Ctx$, whence
\ruleref{EX} returns it to $\Gamma$ at every successor. It is carried forever, and with it the
deadline on $x_\mu$, which is $b$ until the first stop and $b-a$ after. Both are at most $M$, so
$x_\mu$ is bounded along $\rho$ and, by Lemma~\ref{lem:search}(ii), every time-progressing cycle
$\rho$ meets resets it. Only \ruleref{Ustoproll} does, \ruleref{Cov} being blocked while the pursuit
is live, and by \ref{M1} the schedule reaches such a cycle and traverses it without end, so the stop
recurs; its first guard is $x_\mu\in\kI$ and each later one is taken under $x_\mu\le b-a$ measured
from the preceding stop, which are the two bounds. Now let $t\ge t_0$ be a position and
$w_j$ the first stop at or after $t+a$. If $j=1$ then $w_1\le t_0+b\le t+b$; otherwise $w_{j-1}<t+a$,
so $w_j\le w_{j-1}+(b-a)<t+b$. Either way $w_j-t\in\kI$, and $\psi$ holds at $w_j$ while
\ruleref{Ugo} asserts $\varphi$ throughout, which is \eqref{eq:sat_u}.
\end{proof}

\begin{theorem}[Soundness]\label{thm:soundness}
If Algorithm~\ref{alg:otf} returns \SAT{} then $\varphi$ is satisfiable, and $\mathcal A_\varphi$ is
a witness for it.
\end{theorem}
\begin{proof}
By Theorem~\ref{thm:alg} the root lies in $\mathcal M^{*}$, so $\mathcal A_\varphi$ is defined, and
Lemma~\ref{lem:zones}(i) supplies valuations and delays along its paths. It is non-blocking: its only
invariants are the deadlines in $\Deadline$, and cap and stop guard derive from the same interval with the
same strictness, so the stop edge is enabled up to the cap and by \ref{M2} the schedule takes it. By
\ref{M1} the schedule reaches a time-progressing cycle, which Lemma~\ref{lem:zones}(iii) unfolds into
a run whose iterations each elapse a fixed $\delta>0$, so $\mathrm{Runs}^\infty$ is non-empty
everywhere, which is all the semantics asks of it.

We show $\mathcal T,(s_0,\mathbf 0)\models\varphi$ for the induced TTS by induction on $\varphi$
along the marking. Propositional goals hold by saturation and closure of contradictory nodes;
$\EX\psi$ obtains a successor satisfying $\psi$ from its \ruleref{EX} child, and $\AX\psi$ holds at
all successors since each child inherits $\Ctx$. For a timed modality demanded at $N$,
\ruleref{Cov} is conjunctive, so $\mathcal M^{*}$ retains both the discharge child, carrying the
remaining goals, and the launch, which establishes $\mu$ with $x_\mu$ reset, so \kI is measured from
this configuration and the pursuit rides every successor, as a goal under $\mathrm{E}$ and as context
under $\mathrm{A}$. An existential until is witnessed by the path \ref{M2} supplies, extended past
its \ruleref{Ustop} by a divergent run; a universal one cannot be evaded, $\Inv(N)$ blocking time at
the cap while the pursuit is imposed on all successors, and when the demand is persistent this is
Lemma~\ref{lem:recurrence}. \ruleref{Rsplit} discharges a release by partitioning the zone
conjunctively: in the window \ruleref{Runf_in} asserts $\psi$ until $\varphi$ releases it, before it
\ruleref{Runf_before} carries the pursuit on, and past it the obligation has lapsed.
Lemma~\ref{lem:zones}(ii) lifts each conclusion to the positions between nodes.
\end{proof}

\begin{theorem}[Completeness]\label{thm:completeness}
If $\varphi$ is non-punctual and satisfiable then Algorithm~\ref{alg:otf} returns \SAT.
\end{theorem}
\begin{proof}
Completeness goes the other way: the realisable nodes of \S\ref{ss:nodes} are themselves a successful
marking. Writing $R$ for them, we show $R\subseteq F(R)$; then $R\subseteq\mathcal M^{*}$ by
Proposition~\ref{prop:gfp}, and $N_0\in R$ as $\varphi$ is satisfiable, so Theorem~\ref{thm:alg}
returns \SAT.

\ref{D1}: a realisable node has a non-empty zone and a consistent $\Gamma_{\mathsf{prop}}$.
\ref{D2}: let $N\in R$ be realised by $\mathcal T$. All children of an \ANDn-rule are realised by the
same $\mathcal T$: \ruleref{Cov} because $\mathcal T$ satisfies $\Gamma\setminus\{\mu\}$ and $\mu$
alike, \ruleref{Rsplit} because its children partition $Z_N$ and realisability is monotone under
narrowing, \ruleref{Delay} and \ruleref{EX} because $\mathcal T$ admits the passage of time and
supplies a successor for every $\EX\psi$, each inheriting $\Ctx$. At an \ORn-rule $\mathcal T$ selects
a child: a disjunct it satisfies for \ruleref{Disj}; for \ruleref{Uunf} and \ruleref{Uunfroll} the
stop if $\psi$ holds now with $x_\mu\in\kI$ and \ruleref{Ugo} otherwise, the witness being later and
$\varphi$ holding now; for \ruleref{Runf_before} and \ruleref{Runf_in} the stop if $\varphi$ has
released the obligation and the recursive child otherwise. The zone operations being exact, no chosen
zone is empty; and that realisability quantifies universally over $Z_N$ is what lets one child answer
a demand re-imposed at a continuum of onsets.
\ref{D3} and \ref{D4}: fix $N\in R$ and a run $\rho$ of $\mathcal T$ from a configuration for $N$. It
is non-Zeno, so no clock stays both unreset and bounded along it, and following $\rho$ through the
choices above gives a path of $R$ from $N$ which by Lemma~\ref{lem:finiteness} repeats a node; that
segment elapses positive time with every clock reset on it or extrapolated above $M$, hence is a
time-progressing cycle of $R$ reached from $N$. An until pursuit active at $N$ is one $\mathcal T$
fulfils at a finite position of $\rho$, which the path following $\rho$ carries it to.

Non-punctuality enters at the repetition step: a node summarises the pending obligations of a
persistent universal until by $x_\mu$ alone, the time since the last witness capped at $b-a$, and by
Lemma~\ref{lem:recurrence} that summary is faithful exactly when witnesses spaced at most $b-a$ apart
answer every onset, which needs $b>a$.
\end{proof}


\begin{figure}[b]
\centering
\begin{tikzpicture}[x=1mm,y=1mm,font=\scriptsize]
\begin{scope}
\draw[->] (14,22) -- (58,22);
\foreach \x/\lab in {20/1,34/2,48/3} \draw (\x,21.4) -- (\x,22.6) node[above=0pt]{\lab};
\node at (13,17) {$0$};   \draw[line width=1.4pt] (20,17) -- (34,17);
\myvdots{13}{14.7}{0.7}{0.1}
\node at (13,11) {$0.4$}; \draw[line width=1.4pt] (25.6,11) -- (39.6,11);
\myvdots{13}{8.7}{0.7}{0.1}
\node at (13,5) {$0.8$};  \draw[line width=1.4pt] (31.2,5) -- (45.2,5);
\myvdots{13}{2.7}{0.7}{0.1}
\draw[dashed] (33,18.5) -- (33,3.5);
\foreach \y in {17,11,5} \fill (33,\y) circle (0.75);
\node at (36,-1) {(a)\ $\kI=[1,2]$};
\end{scope}
\begin{scope}[xshift=52mm]
\draw[->] (14,22) -- (58,22);
\foreach \x/\lab in {20/1,34/2,48/3} \draw (\x,21.4) -- (\x,22.6) node[above=0pt]{\lab};
\node at (13,17) {$0$};   \fill (20,17) circle (0.75);
\myvdots{13}{14.7}{0.7}{0.1}
\node at (13,11) {$0.4$}; \fill (25.6,11) circle (0.75);
\myvdots{13}{8.7}{0.7}{0.1}
\node at (13,5) {$0.8$};  \fill (31.2,5) circle (0.75);
\myvdots{13}{2.7}{0.7}{0.1}
\node at (36,-1) {(b)\ $\kI=[1,1]$};
\end{scope}
\end{tikzpicture}
\caption{Onsets of a persistent demand and their windows: overlapping for a
non-punctual \kI\ \subref*{fig:non_punctual}, pairwise disjoint for a punctual one
\subref*{fig:punctual}.}
\label{fig:punctual_vs_nonpunctual}%
\phantomsubcaption\label{fig:non_punctual}%
\phantomsubcaption\label{fig:punctual}
\end{figure}

\begin{corollary}\label{cor:depthone}
Satisfiability of non-punctual TCTL is \textsc{ExpTime}-complete, and the procedure is sound and
complete for TCTL of temporal nesting depth at most one, punctual intervals included.
\end{corollary}
\begin{proof}
Membership is Theorems~\ref{thm:alg}--\ref{thm:completeness}; hardness is inherited from CTL
satisfiability~\cite{Clarke99} by fixing every interval to $[0,\infty)$. At depth one every argument
is propositional, so no modality is demanded inside the pursuit of another and no fresh onset arises;
the obligations live at a node are exactly the pursuits it carries, recorded in $\Gamma$, $\Ctx$ and
$\Deadline$ with their elapsed time in $Z$, so a repeat agrees on the full obligation state and
Theorem~\ref{thm:completeness} goes through for every \kI.
\end{proof}

\noindent
Fig.~\ref{fig:punctual_vs_nonpunctual} draws the gap: for $\kI=[a,b]$ with $a<b$ the windows of
nearby onsets overlap, so the pending obligations collapse to the earliest deadline and $x_\mu$
records it, whereas for $\kI=[a,a]$ they are disjoint singletons and what a model owes depends on the
exact offsets of every onset still pending.

\noindent
Depth one covers reachability, safety and inevitability, the queries
\textsc{Uppaal}~\cite{BehrmannLarsenYi2004Uppaal} and \textsc{Romeo}~\cite{lime-TACAS-09} are asked
for; the one common query beyond it, leads-to $\mathrm{AG}(p\rightarrow\mathrm{AF}_\kI q)$, stays
covered by Theorem~\ref{thm:completeness} whenever \kI is non-punctual.


%% =====================================================================
%% APPENDIX A -- deferred proofs.  \input this after \bibliography in
%% ETAPS2026.tex, inside an \appendix block:
%%
%%   \appendix
%%   %% =====================================================================
%% DRAFT / extended version of Section 5.  NOT the file \input by
%% ETAPS2026.tex; the shipped version is sec/section5_OTF.tex.
%%
%% This file is the same section with every proof inline and the
%% punctuality figure kept in place (~6 pages).  The shipped version
%% (~4.6 pages) states Proposition 1 and Lemmas 1-4 without proof and
%% points at the appendix; APPENDIX A below is exactly that appendix,
%% ready to \input after the bibliography.
%%
%% Design notes are at the end of the file.
%% =====================================================================

\section{The Decision Procedure}\label{sec:correct}

Section~\ref{sec:tableau} defines success by a condition on a marking of a graph that no
implementation can afford to build. We make that graph finite and identify the largest successful
marking with a greatest fixpoint (\S\ref{ss:gfp}), give a procedure deciding membership of the root
in it while building only what the search demands (\S\ref{ss:otf}), and prove that procedure sound
for TCTL and complete for the non-punctual fragment (\S\ref{ss:sound}). Throughout $\varphi$ is the
input, $n=|\FL(\varphi)|$, and $M$ the largest constant in its intervals. 

\subsection{The tableau graph and its greatest marking}\label{ss:gfp}

The calculus leaves the order of rule application open; the procedure fixes one. Let
$\FL^{\bullet}(\varphi)$ extend $\FL(\varphi)$ with the pursuit markers $\mu^\bullet,\mu^{\ast}$ and
with $\EX\chi,\AX\chi$ for $\chi$ in it, so that it holds every goal the rules produce and has size
$O(n)$, and fix a linear order on it. Then $\Expand{N}$ is empty if $N$ is closed, and otherwise
fires on the least goal of $\Gamma$ it fits the first applicable rule in the order \ruleref{Conj},
\ruleref{Disj}, \ruleref{AX}, \ruleref{Roll}, \ruleref{Cov}, \ruleref{Rsplit},
\ruleref{Runf_before}, \ruleref{Runf_in}, \ruleref{Uunf}, \ruleref{Uunfroll}; when none applies, $N$
is saturated and $\Expand{N}$ returns the \ruleref{Delay} and \ruleref{EX} children together, both
being conjunctive. Two conventions make this a function of the node: the clock a \ruleref{launch}
allocates is determined by $\mu$ rather than drawn fresh, sound because \ruleref{Cov} blocks a second
launch while a pursuit of $\mu$ is live, so $\clockSet=\{x_\mu\mid\mu$ timed in $\FL(\varphi)\}$ has
size $O(n)$; and acceptance reads clock resets and stop branches off the graph, both properties of an
edge, so each node records the clocks its rule reset and the modality it discharged.

\begin{definition}[Tableau graph]\label{def:graph}
$G(\varphi)$ is the least set of nodes containing $N_0=\tup{\{\varphi\}}{\emptyset}{\{\mathbf
0\}}{\emptyset}{\emptyset}$ and closed under $\Expand{}$, two nodes being identified when they agree
componentwise and their zones agree modulo $\cong_M$, with an edge $N\to N'$ for each
$N'\in\Expand{N}$.
\end{definition}

\begin{lemma}[Finiteness]\label{lem:finiteness}
$|G(\varphi)|\le 2^{O(n)}\cdot M^{O(n^{2})}$, and every node has $O(n)$ children.
\end{lemma}
\begin{proof}
$\Gamma$, $\Ctx$, the reset set and the discharge marker range over subsets of
$\FL^{\bullet}(\varphi)$ or of $\clockSet$; $\obligations$ is fixed by its domain; and $\Deadline$ gives each
$x_\mu$ one of the bounds $\sup\kI$, $\sup\kI-\inf\kI$ with a strictness flag. Each contributes
$2^{O(n)}$. A canonical $\cong_M$-zone over $K\le n+1$ clocks is a $(K{+}1)^{2}$ DBM with entries
$\infty$ or in $\{-M,\dots,M\}\times\{<,\le\}$, giving $(4M{+}3)^{(n+2)^{2}}=M^{O(n^{2})}$ zones.
Every rule has at most three children except \ruleref{EX}, with one per $\EX\psi\in\Gamma$.
\end{proof}

\noindent
Identification up to $\cong_M$ dispenses with a termination argument for expansion, and so lets the
calculus afford rules that re-impose their own premise: a chain of applications that would regress
forever revisits a node and closes a cycle, which crosses no \ruleref{Delay} step and is discarded by
\ref{M1}.

A marking contains no closed node, at least one child of each of its \ORn-nodes and all children of
each of its \ANDn-nodes, and is successful when every member satisfies \ref{M1} and \ref{M2}. Both
clauses quantify over paths \emph{inside} the marking, so a successful marking is a set that survives
the deletions it prescribes. For $A\subseteq G(\varphi)$ let $F(A)$ be $A$ minus:
\begin{enumerate}[start=1,label={D\arabic*)}, itemindent=1.6em, ref=(D\arabic*), itemsep=0.02em]
\item\label{D1} the closed nodes;
\item\label{D2} the \ORn-nodes with no child in $A$, and the \ANDn-nodes with a child outside $A$;
\item\label{D3} the nodes with no $A$-path to a node on a time-progressing cycle of $A$;
\item\label{D4} for each until $\mu$, the nodes pursuing $\mu$ from which no $A$-path carrying the
      pursuit throughout reaches a stop for $\mu$ taken where the zone meets $x_\mu\in\kI$.
\end{enumerate}

\begin{proposition}\label{prop:gfp}
$F$ is monotone, and $A$ is a successful marking iff $A\subseteq F(A)$. Hence the successful markings
have a greatest element $\mathcal M^{*}=\nu F$, and the tableau succeeds on $\varphi$ iff
$N_0\in\mathcal M^{*}$.
\end{proposition}
\begin{proof}
\ref{D1} ignores $A$; a node kept by \ref{D2} has its witnessing child, or all its children, in $A$
and hence in any superset; \ref{D3} and \ref{D4} ask for a path in $A$, which survives enlargement.
So $F$ is monotone on the finite lattice $2^{G(\varphi)}$, and $A\subseteq F(A)$ unfolds to the
clauses defining a successful marking. Knaster--Tarski gives a greatest post-fixpoint, equal to $\nu
F$ and to the union of all post-fixpoints, and the tableau succeeds iff one of them contains $N_0$.
\end{proof}

\subsection{The on-the-fly procedure}\label{ss:otf}

Algorithm~\ref{alg:otf} decides $N_0\in\mathcal M^{*}$ without ever holding $G(\varphi)$.
$\Expand{}$ applies the rules the first time the search reaches a node and returns the recorded
children of a $\cong_M$-equal node otherwise, so the graph exists only as far as the search has
demanded it. Tarjan's indices, used as a black box, ride along; a back edge yields a lowlink update
and nothing else, a node still on the stack having no verdict and none being guessed for it. The
third line of the loop is where the search stops at the first witness: a child settles its parent
when the parent is an \ORn-node and the child holds, or an \ANDn-node and it fails, which is what
$\val[w]=[\,v$ is $\ORn\,]$ says, and the parent retires with its remaining children unbuilt; the guard $\mathrm{top}(\Stack)=v$ makes this well defined, an earlier child
possibly having left members of $v$'s component above $v$ on the stack. A component is decided the
moment it closes, against the verdicts of the children it leaves: writing
$\Out_C=\{w\notin C\mid w\in\ch[v]$ for some $v\in C$, $\val[w]=\top\}$ for those already accepted,
$\nu F_C$ is the largest $A\subseteq C$ with $A\subseteq F(A\cup\Out_C)$, computed by deletion from
$A=C$, each round applying \ref{D1} and \ref{D2} pointwise and then the backward reachabilities
\ref{D3}, \ref{D4} over the survivors. A deletion can disconnect $A$, so the decomposition is
recomputed each round.

\begin{algorithm}[!tb]
\footnotesize
\DontPrintSemicolon
\SetAlgoLined
\SetKwProg{Fn}{procedure}{}{end}

\Fn{\Decide{$\varphi$}}{
  $\val[v]\gets{?}$ and $\idx[v]\gets\bot$ for all $v$; \ $\Stack\gets\varepsilon$; \ $\cnt\gets0$\;
  \Explore{$N_0$} \ \textbf{where} $N_0=\tup{\{\varphi\}}{\emptyset}{\{\mathbf
    0\}}{\emptyset}{\emptyset}$; \quad \Return{\SAT{} \textnormal{iff} $\val[N_0]=\top$}\;
}
\BlankLine
\Fn{\Explore{$v$}}{
  $\idx[v]\gets\low[v]\gets\cnt{+}{+}$; \ $\mathrm{push}(\Stack,v)$; \
  $\ch[v]\gets\Expand{v}$ \tcp*{memoised on $\cong_M$}
  \ForEach{$w\in\ch[v]$}{
    \lIf{$\idx[w]=\bot$}{\Explore{$w$}; \ $\low[v]\gets\min(\low[v],\low[w])$}
    \lElseIf{$w\in\Stack$}{$\low[v]\gets\min(\low[v],\idx[w])$}
    \lIf{$\val[w]=[\,v\text{ is }\ORn\,]$ \textnormal{\textbf{and}} $\mathrm{top}(\Stack)=v$}
        {$\val[v]\gets\val[w]$; \ $\mathrm{pop}(\Stack)$; \Return}
  }
  \lIf{$\low[v]=\idx[v]$}{$C\gets$ pop $\Stack$ down to $v$; \
    $\val[u]\gets[\,u\in\nu F_C\,]$ for $u\in C$}
}
\caption{On-the-fly satisfiability.}\label{alg:otf}
\end{algorithm}

\begin{lemma}\label{lem:search}
\emph{(i)} If $\val[w]$ is final and settles $v$, then $v\in\mathcal M^{*}$ iff $w\in\mathcal M^{*}$,
and leaving the remaining children of $v$ unexpanded changes $\mathcal M^{*}$ at no node.
\emph{(ii)} A strongly connected $T$ with an internal edge carries a time-progressing cycle iff some
node of $T$ is a \ruleref{Delay}-successor and every clock is reset somewhere in $T$ or unbounded
everywhere in $T$.
\emph{(iii)} If $C$ is a maximal strongly connected set with $\val[c]=[\,c\in\mathcal M^{*}\,]$ for
every child $c\notin C$ of a node of $C$, then $\mathcal M^{*}\cap C=\nu F_C$.
\end{lemma}
\begin{proof}
(i) The first claim is \ref{D2}. The explored subgraph is closed under $\Expand{}$ except at retired
nodes, so a path of $\mathcal M^{*}$ using an unexpanded child of $v$ passes through $v$; if
$v\notin\mathcal M^{*}$ there is no such path, and if $v\in\mathcal M^{*}$ then $v$ satisfies
\ref{M1}, and \ref{M2} for each pursuit active at it, so a node reaching $v$ inherits both witnesses
by prefixing the path to $v$, which carries a pursuit exactly when it is active at $v$.

(ii) $(\Leftarrow)$ Concatenating paths in $T$ gives a cycle through the \ruleref{Delay}-successor,
which therefore elapses positive time, and through one reset node of every clock $T$ resets; the rest
are unbounded throughout. $(\Rightarrow)$ A time-progressing cycle in $T$ elapses positive time,
hence crosses a \ruleref{Delay} step; and a clock reset nowhere in $T$ is forced unbounded on that
cycle, whence unbounded on all of $T$, since a path inside $T$ never resets it and its value, already
above $M$, only grows.

(iii) $C$ being maximal, a path of $\mathcal M^{*}$ from $N\in C$ either stays in $C$, and a witness
for \ref{M1} or \ref{M2} then lies inside $C$; or it leaves $C$ once and for all at some
$c\in\mathcal M^{*}\setminus C$, which by hypothesis lies in $\Out_C$ and itself satisfies \ref{M1},
and \ref{M2} for every pursuit active at it, so prefixing the path from $N$ to $c$ turns either
witness into one for $N$. Hence $\mathcal M^{*}\cap C\subseteq F(\mathcal M^{*}\cap C\cup\Out_C)$.
Conversely any $A\subseteq C$ with $A\subseteq F(A\cup\Out_C)$ extends by the nodes of
$\mathcal M^{*}$ reachable from $\Out_C$ to a post-fixpoint of $F$, which Proposition~\ref{prop:gfp}
places inside $\mathcal M^{*}$. The greatest elements coincide.
\end{proof}

\begin{theorem}[Adequacy and complexity]\label{thm:alg}
Algorithm~\ref{alg:otf} terminates and retires every node it builds with
$\val[N]=[\,N\in\mathcal M^{*}\,]$; it returns \SAT{} iff $N_0\in\mathcal M^{*}$, and runs in time
$2^{\mathrm{poly}(n)}$.
\end{theorem}
\begin{proof}
By Lemma~\ref{lem:finiteness} with memoisation the traversal is a depth-first search of a finite
graph, so it terminates. Components are popped in reverse topological order, and by the Tarjan
invariant every child of $C$ outside $C$ has retired when $C$ is popped, one still on the stack
dragging the lowlink of $C$'s root below its index; so $\Out_C$ is final and every pop meets
Lemma~\ref{lem:search}(iii). The deletion loop decreases on a finite set and each step instantiates a
clause of $F$, exact by Lemma~\ref{lem:search}(ii), so it reaches $\nu F_C=\mathcal M^{*}\cap C$, and
by Lemma~\ref{lem:search}(i) pruning preserves every verdict.

For the bound let $m$ be that of Lemma~\ref{lem:finiteness}. As $M$ occupies at most $n$ input
symbols, $\log m=O(n)+O(n^{2})\log M=\mathrm{poly}(n)$. The explored graph has at most $m$ nodes and
$O(nm)$ edges, and memoisation on $\cong_M$-classes expands each once, avoiding the doubly
exponential blow-up of an ancestor-only unfolding. A component takes at most $|C|$ rounds,
each deleting a node or stopping, of one decomposition and $O(n)$ backward reachabilities linear in
the subgraph, and the components partition the nodes, so the run costs $2^{\mathrm{poly}(n)}$.
\end{proof}

\subsection{Soundness and Completeness}\label{ss:sound}

\begin{lemma}\label{lem:zones}
Let $N_0,\dots,N_k$ be a path of $G(\varphi)$ with non-empty zones. \emph{(i)} For every $\nu_k\in
Z_k$ there are $\nu_i\in Z_i$ and delays realising every step, positive at every \ruleref{Delay}
step. \emph{(ii)} If a delay crosses no integer hyperplane $x=b$ with $b\le M$, every subformula of
$\varphi$ has constant truth on the open segment it sweeps. \emph{(iii)} A cycle of $G(\varphi)$ has
a non-Zeno unfolding iff it is time-progressing, and then one whose iterations elapse a constant
positive amount.
\end{lemma}
\begin{proof}
(i) Every zone operation a rule performs (intersection with a single-clock guard, reset,
time-successor, $\cong_M$-canonicalisation) back-propagates, non-empty precisely because the
successor valuation lies in its image~\cite{Bouyer05TimedAutomata}; thread backwards from $\nu_k$,
noting that $\Dopen Z$ and $\Dbd Z$ lie strictly ahead of $Z$ in time. (ii) No clock crosses an
integer at or below $M$, so no integer part changes and no fractional ordering wraps, while clocks
above $M$ are identified by extrapolation; the segment lies in one region, on which TCTL truth is
constant~\cite{AlurCourcoubetisDill1990,Bouyer05TimedAutomata}, and its endpoints are nodes of their
own by the integer split of \ruleref{Delay}. (iii) A cycle elapsing no time takes zero time over
infinitely many traversals; otherwise some clock is neither reset nor extrapolated above $M$, so it
keeps a bound $c\le M$, grows by exactly the delay elapsed, and caps the total at $c$. Conversely (i)
yields a traversal elapsing some $\delta>0$, and the cycle returning to the same node, no rule
distinguishes two visits and no budget depletes, so $\delta$ recurs at every iteration.
\end{proof}

\noindent
\textbf{From a marking to an automaton.} Call a node of $\mathcal M^{*}$ a \emph{state} if it is
saturated, so that its children are the successor steps; the remaining, \emph{expansion} edges
neither elapse time nor change location, and reset a clock only through a \ruleref{launch} or a
\ruleref{Ustoproll}. At most $n+1$ obligations are pending at $N$, the divergence requirement of
\ref{M1} and the until pursuits active there, and \ref{M1}, \ref{M2} give a finite path $\pi_{N,o}$
of $\mathcal M^{*}$ discharging each. A \emph{schedule} selects at every \ORn-node the child
$\pi_{N,o}$ prescribes and moves round-robin to the next pending obligation once $o$ is discharged,
so its memory, the pair (obligation, position on $\pi_{N,o}$), is finite. The witness
$\mathcal A_\varphi$ has as locations the pairs (state, schedule memory), labelled by
$\Gamma_{\mathsf{prop}}$ and carrying $\Inv(N)$, and as edges the scheduled successor steps, guarded
by the zone of the next state reached through expansion edges and resetting the clocks that segment
resets.

\begin{lemma}[Recurrence]\label{lem:recurrence}
Let $\mu=\mathrm{A}(\varphi\,\mathbf{U}_\kI\,\psi)$ with $a=\inf\kI$, $b=\sup\kI<\infty$, and let
$\rho$ be a run of $\mathcal A_\varphi$ through a node carrying $\mu^{\ast}$, launched at $t_0$. Then
\ruleref{Ustoproll} fires on $\rho$ at times $t_0<w_1<w_2<\cdots$ with $w_1-t_0\in\kI$ and
$w_{j+1}-w_j\le b-a$, and $\rho\models\varphi\,\mathbf{U}_\kI\,\psi$ from every position at or after
$t_0$.
\end{lemma}
\begin{proof}
No rule drops $\mu^{\ast}$: \ruleref{Ustoproll} keeps it and \ruleref{Ugo} moves it to $\Ctx$, whence
\ruleref{EX} returns it to $\Gamma$ at every successor. It is carried forever, and with it the
deadline on $x_\mu$, which is $b$ until the first stop and $b-a$ after. Both are at most $M$, so
$x_\mu$ is bounded along $\rho$ and, by Lemma~\ref{lem:search}(ii), every time-progressing cycle
$\rho$ meets resets it. Only \ruleref{Ustoproll} does, \ruleref{Cov} being blocked while the pursuit
is live, and by \ref{M1} the schedule reaches such a cycle and traverses it without end, so the stop
recurs; its first guard is $x_\mu\in\kI$ and each later one is taken under $x_\mu\le b-a$ measured
from the preceding stop, which are the two bounds. Now let $t\ge t_0$ be a position and
$w_j$ the first stop at or after $t+a$. If $j=1$ then $w_1\le t_0+b\le t+b$; otherwise $w_{j-1}<t+a$,
so $w_j\le w_{j-1}+(b-a)<t+b$. Either way $w_j-t\in\kI$, and $\psi$ holds at $w_j$ while
\ruleref{Ugo} asserts $\varphi$ throughout, which is \eqref{eq:sat_u}.
\end{proof}

\begin{theorem}[Soundness]\label{thm:soundness}
If Algorithm~\ref{alg:otf} returns \SAT{} then $\varphi$ is satisfiable, and $\mathcal A_\varphi$ is
a witness for it.
\end{theorem}
\begin{proof}
By Theorem~\ref{thm:alg} the root lies in $\mathcal M^{*}$, so $\mathcal A_\varphi$ is defined, and
Lemma~\ref{lem:zones}(i) supplies valuations and delays along its paths. It is non-blocking: its only
invariants are the deadlines in $\Deadline$, and cap and stop guard derive from the same interval with the
same strictness, so the stop edge is enabled up to the cap and by \ref{M2} the schedule takes it. By
\ref{M1} the schedule reaches a time-progressing cycle, which Lemma~\ref{lem:zones}(iii) unfolds into
a run whose iterations each elapse a fixed $\delta>0$, so $\mathrm{Runs}^\infty$ is non-empty
everywhere, which is all the semantics asks of it.

We show $\mathcal T,(s_0,\mathbf 0)\models\varphi$ for the induced TTS by induction on $\varphi$
along the marking. Propositional goals hold by saturation and closure of contradictory nodes;
$\EX\psi$ obtains a successor satisfying $\psi$ from its \ruleref{EX} child, and $\AX\psi$ holds at
all successors since each child inherits $\Ctx$. For a timed modality demanded at $N$,
\ruleref{Cov} is conjunctive, so $\mathcal M^{*}$ retains both the discharge child, carrying the
remaining goals, and the launch, which establishes $\mu$ with $x_\mu$ reset, so \kI is measured from
this configuration and the pursuit rides every successor, as a goal under $\mathrm{E}$ and as context
under $\mathrm{A}$. An existential until is witnessed by the path \ref{M2} supplies, extended past
its \ruleref{Ustop} by a divergent run; a universal one cannot be evaded, $\Inv(N)$ blocking time at
the cap while the pursuit is imposed on all successors, and when the demand is persistent this is
Lemma~\ref{lem:recurrence}. \ruleref{Rsplit} discharges a release by partitioning the zone
conjunctively: in the window \ruleref{Runf_in} asserts $\psi$ until $\varphi$ releases it, before it
\ruleref{Runf_before} carries the pursuit on, and past it the obligation has lapsed.
Lemma~\ref{lem:zones}(ii) lifts each conclusion to the positions between nodes.
\end{proof}

\begin{theorem}[Completeness]\label{thm:completeness}
If $\varphi$ is non-punctual and satisfiable then Algorithm~\ref{alg:otf} returns \SAT.
\end{theorem}
\begin{proof}
Completeness goes the other way: the realisable nodes of \S\ref{ss:nodes} are themselves a successful
marking. Writing $R$ for them, we show $R\subseteq F(R)$; then $R\subseteq\mathcal M^{*}$ by
Proposition~\ref{prop:gfp}, and $N_0\in R$ as $\varphi$ is satisfiable, so Theorem~\ref{thm:alg}
returns \SAT.

\ref{D1}: a realisable node has a non-empty zone and a consistent $\Gamma_{\mathsf{prop}}$.
\ref{D2}: let $N\in R$ be realised by $\mathcal T$. All children of an \ANDn-rule are realised by the
same $\mathcal T$: \ruleref{Cov} because $\mathcal T$ satisfies $\Gamma\setminus\{\mu\}$ and $\mu$
alike, \ruleref{Rsplit} because its children partition $Z_N$ and realisability is monotone under
narrowing, \ruleref{Delay} and \ruleref{EX} because $\mathcal T$ admits the passage of time and
supplies a successor for every $\EX\psi$, each inheriting $\Ctx$. At an \ORn-rule $\mathcal T$ selects
a child: a disjunct it satisfies for \ruleref{Disj}; for \ruleref{Uunf} and \ruleref{Uunfroll} the
stop if $\psi$ holds now with $x_\mu\in\kI$ and \ruleref{Ugo} otherwise, the witness being later and
$\varphi$ holding now; for \ruleref{Runf_before} and \ruleref{Runf_in} the stop if $\varphi$ has
released the obligation and the recursive child otherwise. The zone operations being exact, no chosen
zone is empty; and that realisability quantifies universally over $Z_N$ is what lets one child answer
a demand re-imposed at a continuum of onsets.
\ref{D3} and \ref{D4}: fix $N\in R$ and a run $\rho$ of $\mathcal T$ from a configuration for $N$. It
is non-Zeno, so no clock stays both unreset and bounded along it, and following $\rho$ through the
choices above gives a path of $R$ from $N$ which by Lemma~\ref{lem:finiteness} repeats a node; that
segment elapses positive time with every clock reset on it or extrapolated above $M$, hence is a
time-progressing cycle of $R$ reached from $N$. An until pursuit active at $N$ is one $\mathcal T$
fulfils at a finite position of $\rho$, which the path following $\rho$ carries it to.

Non-punctuality enters at the repetition step: a node summarises the pending obligations of a
persistent universal until by $x_\mu$ alone, the time since the last witness capped at $b-a$, and by
Lemma~\ref{lem:recurrence} that summary is faithful exactly when witnesses spaced at most $b-a$ apart
answer every onset, which needs $b>a$.
\end{proof}


\begin{figure}[b]
\centering
\begin{tikzpicture}[x=1mm,y=1mm,font=\scriptsize]
\begin{scope}
\draw[->] (14,22) -- (58,22);
\foreach \x/\lab in {20/1,34/2,48/3} \draw (\x,21.4) -- (\x,22.6) node[above=0pt]{\lab};
\node at (13,17) {$0$};   \draw[line width=1.4pt] (20,17) -- (34,17);
\myvdots{13}{14.7}{0.7}{0.1}
\node at (13,11) {$0.4$}; \draw[line width=1.4pt] (25.6,11) -- (39.6,11);
\myvdots{13}{8.7}{0.7}{0.1}
\node at (13,5) {$0.8$};  \draw[line width=1.4pt] (31.2,5) -- (45.2,5);
\myvdots{13}{2.7}{0.7}{0.1}
\draw[dashed] (33,18.5) -- (33,3.5);
\foreach \y in {17,11,5} \fill (33,\y) circle (0.75);
\node at (36,-1) {(a)\ $\kI=[1,2]$};
\end{scope}
\begin{scope}[xshift=52mm]
\draw[->] (14,22) -- (58,22);
\foreach \x/\lab in {20/1,34/2,48/3} \draw (\x,21.4) -- (\x,22.6) node[above=0pt]{\lab};
\node at (13,17) {$0$};   \fill (20,17) circle (0.75);
\myvdots{13}{14.7}{0.7}{0.1}
\node at (13,11) {$0.4$}; \fill (25.6,11) circle (0.75);
\myvdots{13}{8.7}{0.7}{0.1}
\node at (13,5) {$0.8$};  \fill (31.2,5) circle (0.75);
\myvdots{13}{2.7}{0.7}{0.1}
\node at (36,-1) {(b)\ $\kI=[1,1]$};
\end{scope}
\end{tikzpicture}
\caption{Onsets of a persistent demand and their windows: overlapping for a
non-punctual \kI\ \subref*{fig:non_punctual}, pairwise disjoint for a punctual one
\subref*{fig:punctual}.}
\label{fig:punctual_vs_nonpunctual}%
\phantomsubcaption\label{fig:non_punctual}%
\phantomsubcaption\label{fig:punctual}
\end{figure}

\begin{corollary}\label{cor:depthone}
Satisfiability of non-punctual TCTL is \textsc{ExpTime}-complete, and the procedure is sound and
complete for TCTL of temporal nesting depth at most one, punctual intervals included.
\end{corollary}
\begin{proof}
Membership is Theorems~\ref{thm:alg}--\ref{thm:completeness}; hardness is inherited from CTL
satisfiability~\cite{Clarke99} by fixing every interval to $[0,\infty)$. At depth one every argument
is propositional, so no modality is demanded inside the pursuit of another and no fresh onset arises;
the obligations live at a node are exactly the pursuits it carries, recorded in $\Gamma$, $\Ctx$ and
$\Deadline$ with their elapsed time in $Z$, so a repeat agrees on the full obligation state and
Theorem~\ref{thm:completeness} goes through for every \kI.
\end{proof}

\noindent
Fig.~\ref{fig:punctual_vs_nonpunctual} draws the gap: for $\kI=[a,b]$ with $a<b$ the windows of
nearby onsets overlap, so the pending obligations collapse to the earliest deadline and $x_\mu$
records it, whereas for $\kI=[a,a]$ they are disjoint singletons and what a model owes depends on the
exact offsets of every onset still pending.

\noindent
Depth one covers reachability, safety and inevitability, the queries
\textsc{Uppaal}~\cite{BehrmannLarsenYi2004Uppaal} and \textsc{Romeo}~\cite{lime-TACAS-09} are asked
for; the one common query beyond it, leads-to $\mathrm{AG}(p\rightarrow\mathrm{AF}_\kI q)$, stays
covered by Theorem~\ref{thm:completeness} whenever \kI is non-punctual.


%% =====================================================================
%% APPENDIX A -- deferred proofs.  \input this after \bibliography in
%% ETAPS2026.tex, inside an \appendix block:
%%
%%   \appendix
%%   %% =====================================================================
%% DRAFT / extended version of Section 5.  NOT the file \input by
%% ETAPS2026.tex; the shipped version is sec/section5_OTF.tex.
%%
%% This file is the same section with every proof inline and the
%% punctuality figure kept in place (~6 pages).  The shipped version
%% (~4.6 pages) states Proposition 1 and Lemmas 1-4 without proof and
%% points at the appendix; APPENDIX A below is exactly that appendix,
%% ready to \input after the bibliography.
%%
%% Design notes are at the end of the file.
%% =====================================================================

\section{The Decision Procedure}\label{sec:correct}

Section~\ref{sec:tableau} defines success by a condition on a marking of a graph that no
implementation can afford to build. We make that graph finite and identify the largest successful
marking with a greatest fixpoint (\S\ref{ss:gfp}), give a procedure deciding membership of the root
in it while building only what the search demands (\S\ref{ss:otf}), and prove that procedure sound
for TCTL and complete for the non-punctual fragment (\S\ref{ss:sound}). Throughout $\varphi$ is the
input, $n=|\FL(\varphi)|$, and $M$ the largest constant in its intervals. 

\subsection{The tableau graph and its greatest marking}\label{ss:gfp}

The calculus leaves the order of rule application open; the procedure fixes one. Let
$\FL^{\bullet}(\varphi)$ extend $\FL(\varphi)$ with the pursuit markers $\mu^\bullet,\mu^{\ast}$ and
with $\EX\chi,\AX\chi$ for $\chi$ in it, so that it holds every goal the rules produce and has size
$O(n)$, and fix a linear order on it. Then $\Expand{N}$ is empty if $N$ is closed, and otherwise
fires on the least goal of $\Gamma$ it fits the first applicable rule in the order \ruleref{Conj},
\ruleref{Disj}, \ruleref{AX}, \ruleref{Roll}, \ruleref{Cov}, \ruleref{Rsplit},
\ruleref{Runf_before}, \ruleref{Runf_in}, \ruleref{Uunf}, \ruleref{Uunfroll}; when none applies, $N$
is saturated and $\Expand{N}$ returns the \ruleref{Delay} and \ruleref{EX} children together, both
being conjunctive. Two conventions make this a function of the node: the clock a \ruleref{launch}
allocates is determined by $\mu$ rather than drawn fresh, sound because \ruleref{Cov} blocks a second
launch while a pursuit of $\mu$ is live, so $\clockSet=\{x_\mu\mid\mu$ timed in $\FL(\varphi)\}$ has
size $O(n)$; and acceptance reads clock resets and stop branches off the graph, both properties of an
edge, so each node records the clocks its rule reset and the modality it discharged.

\begin{definition}[Tableau graph]\label{def:graph}
$G(\varphi)$ is the least set of nodes containing $N_0=\tup{\{\varphi\}}{\emptyset}{\{\mathbf
0\}}{\emptyset}{\emptyset}$ and closed under $\Expand{}$, two nodes being identified when they agree
componentwise and their zones agree modulo $\cong_M$, with an edge $N\to N'$ for each
$N'\in\Expand{N}$.
\end{definition}

\begin{lemma}[Finiteness]\label{lem:finiteness}
$|G(\varphi)|\le 2^{O(n)}\cdot M^{O(n^{2})}$, and every node has $O(n)$ children.
\end{lemma}
\begin{proof}
$\Gamma$, $\Ctx$, the reset set and the discharge marker range over subsets of
$\FL^{\bullet}(\varphi)$ or of $\clockSet$; $\obligations$ is fixed by its domain; and $\Deadline$ gives each
$x_\mu$ one of the bounds $\sup\kI$, $\sup\kI-\inf\kI$ with a strictness flag. Each contributes
$2^{O(n)}$. A canonical $\cong_M$-zone over $K\le n+1$ clocks is a $(K{+}1)^{2}$ DBM with entries
$\infty$ or in $\{-M,\dots,M\}\times\{<,\le\}$, giving $(4M{+}3)^{(n+2)^{2}}=M^{O(n^{2})}$ zones.
Every rule has at most three children except \ruleref{EX}, with one per $\EX\psi\in\Gamma$.
\end{proof}

\noindent
Identification up to $\cong_M$ dispenses with a termination argument for expansion, and so lets the
calculus afford rules that re-impose their own premise: a chain of applications that would regress
forever revisits a node and closes a cycle, which crosses no \ruleref{Delay} step and is discarded by
\ref{M1}.

A marking contains no closed node, at least one child of each of its \ORn-nodes and all children of
each of its \ANDn-nodes, and is successful when every member satisfies \ref{M1} and \ref{M2}. Both
clauses quantify over paths \emph{inside} the marking, so a successful marking is a set that survives
the deletions it prescribes. For $A\subseteq G(\varphi)$ let $F(A)$ be $A$ minus:
\begin{enumerate}[start=1,label={D\arabic*)}, itemindent=1.6em, ref=(D\arabic*), itemsep=0.02em]
\item\label{D1} the closed nodes;
\item\label{D2} the \ORn-nodes with no child in $A$, and the \ANDn-nodes with a child outside $A$;
\item\label{D3} the nodes with no $A$-path to a node on a time-progressing cycle of $A$;
\item\label{D4} for each until $\mu$, the nodes pursuing $\mu$ from which no $A$-path carrying the
      pursuit throughout reaches a stop for $\mu$ taken where the zone meets $x_\mu\in\kI$.
\end{enumerate}

\begin{proposition}\label{prop:gfp}
$F$ is monotone, and $A$ is a successful marking iff $A\subseteq F(A)$. Hence the successful markings
have a greatest element $\mathcal M^{*}=\nu F$, and the tableau succeeds on $\varphi$ iff
$N_0\in\mathcal M^{*}$.
\end{proposition}
\begin{proof}
\ref{D1} ignores $A$; a node kept by \ref{D2} has its witnessing child, or all its children, in $A$
and hence in any superset; \ref{D3} and \ref{D4} ask for a path in $A$, which survives enlargement.
So $F$ is monotone on the finite lattice $2^{G(\varphi)}$, and $A\subseteq F(A)$ unfolds to the
clauses defining a successful marking. Knaster--Tarski gives a greatest post-fixpoint, equal to $\nu
F$ and to the union of all post-fixpoints, and the tableau succeeds iff one of them contains $N_0$.
\end{proof}

\subsection{The on-the-fly procedure}\label{ss:otf}

Algorithm~\ref{alg:otf} decides $N_0\in\mathcal M^{*}$ without ever holding $G(\varphi)$.
$\Expand{}$ applies the rules the first time the search reaches a node and returns the recorded
children of a $\cong_M$-equal node otherwise, so the graph exists only as far as the search has
demanded it. Tarjan's indices, used as a black box, ride along; a back edge yields a lowlink update
and nothing else, a node still on the stack having no verdict and none being guessed for it. The
third line of the loop is where the search stops at the first witness: a child settles its parent
when the parent is an \ORn-node and the child holds, or an \ANDn-node and it fails, which is what
$\val[w]=[\,v$ is $\ORn\,]$ says, and the parent retires with its remaining children unbuilt; the guard $\mathrm{top}(\Stack)=v$ makes this well defined, an earlier child
possibly having left members of $v$'s component above $v$ on the stack. A component is decided the
moment it closes, against the verdicts of the children it leaves: writing
$\Out_C=\{w\notin C\mid w\in\ch[v]$ for some $v\in C$, $\val[w]=\top\}$ for those already accepted,
$\nu F_C$ is the largest $A\subseteq C$ with $A\subseteq F(A\cup\Out_C)$, computed by deletion from
$A=C$, each round applying \ref{D1} and \ref{D2} pointwise and then the backward reachabilities
\ref{D3}, \ref{D4} over the survivors. A deletion can disconnect $A$, so the decomposition is
recomputed each round.

\begin{algorithm}[!tb]
\footnotesize
\DontPrintSemicolon
\SetAlgoLined
\SetKwProg{Fn}{procedure}{}{end}

\Fn{\Decide{$\varphi$}}{
  $\val[v]\gets{?}$ and $\idx[v]\gets\bot$ for all $v$; \ $\Stack\gets\varepsilon$; \ $\cnt\gets0$\;
  \Explore{$N_0$} \ \textbf{where} $N_0=\tup{\{\varphi\}}{\emptyset}{\{\mathbf
    0\}}{\emptyset}{\emptyset}$; \quad \Return{\SAT{} \textnormal{iff} $\val[N_0]=\top$}\;
}
\BlankLine
\Fn{\Explore{$v$}}{
  $\idx[v]\gets\low[v]\gets\cnt{+}{+}$; \ $\mathrm{push}(\Stack,v)$; \
  $\ch[v]\gets\Expand{v}$ \tcp*{memoised on $\cong_M$}
  \ForEach{$w\in\ch[v]$}{
    \lIf{$\idx[w]=\bot$}{\Explore{$w$}; \ $\low[v]\gets\min(\low[v],\low[w])$}
    \lElseIf{$w\in\Stack$}{$\low[v]\gets\min(\low[v],\idx[w])$}
    \lIf{$\val[w]=[\,v\text{ is }\ORn\,]$ \textnormal{\textbf{and}} $\mathrm{top}(\Stack)=v$}
        {$\val[v]\gets\val[w]$; \ $\mathrm{pop}(\Stack)$; \Return}
  }
  \lIf{$\low[v]=\idx[v]$}{$C\gets$ pop $\Stack$ down to $v$; \
    $\val[u]\gets[\,u\in\nu F_C\,]$ for $u\in C$}
}
\caption{On-the-fly satisfiability.}\label{alg:otf}
\end{algorithm}

\begin{lemma}\label{lem:search}
\emph{(i)} If $\val[w]$ is final and settles $v$, then $v\in\mathcal M^{*}$ iff $w\in\mathcal M^{*}$,
and leaving the remaining children of $v$ unexpanded changes $\mathcal M^{*}$ at no node.
\emph{(ii)} A strongly connected $T$ with an internal edge carries a time-progressing cycle iff some
node of $T$ is a \ruleref{Delay}-successor and every clock is reset somewhere in $T$ or unbounded
everywhere in $T$.
\emph{(iii)} If $C$ is a maximal strongly connected set with $\val[c]=[\,c\in\mathcal M^{*}\,]$ for
every child $c\notin C$ of a node of $C$, then $\mathcal M^{*}\cap C=\nu F_C$.
\end{lemma}
\begin{proof}
(i) The first claim is \ref{D2}. The explored subgraph is closed under $\Expand{}$ except at retired
nodes, so a path of $\mathcal M^{*}$ using an unexpanded child of $v$ passes through $v$; if
$v\notin\mathcal M^{*}$ there is no such path, and if $v\in\mathcal M^{*}$ then $v$ satisfies
\ref{M1}, and \ref{M2} for each pursuit active at it, so a node reaching $v$ inherits both witnesses
by prefixing the path to $v$, which carries a pursuit exactly when it is active at $v$.

(ii) $(\Leftarrow)$ Concatenating paths in $T$ gives a cycle through the \ruleref{Delay}-successor,
which therefore elapses positive time, and through one reset node of every clock $T$ resets; the rest
are unbounded throughout. $(\Rightarrow)$ A time-progressing cycle in $T$ elapses positive time,
hence crosses a \ruleref{Delay} step; and a clock reset nowhere in $T$ is forced unbounded on that
cycle, whence unbounded on all of $T$, since a path inside $T$ never resets it and its value, already
above $M$, only grows.

(iii) $C$ being maximal, a path of $\mathcal M^{*}$ from $N\in C$ either stays in $C$, and a witness
for \ref{M1} or \ref{M2} then lies inside $C$; or it leaves $C$ once and for all at some
$c\in\mathcal M^{*}\setminus C$, which by hypothesis lies in $\Out_C$ and itself satisfies \ref{M1},
and \ref{M2} for every pursuit active at it, so prefixing the path from $N$ to $c$ turns either
witness into one for $N$. Hence $\mathcal M^{*}\cap C\subseteq F(\mathcal M^{*}\cap C\cup\Out_C)$.
Conversely any $A\subseteq C$ with $A\subseteq F(A\cup\Out_C)$ extends by the nodes of
$\mathcal M^{*}$ reachable from $\Out_C$ to a post-fixpoint of $F$, which Proposition~\ref{prop:gfp}
places inside $\mathcal M^{*}$. The greatest elements coincide.
\end{proof}

\begin{theorem}[Adequacy and complexity]\label{thm:alg}
Algorithm~\ref{alg:otf} terminates and retires every node it builds with
$\val[N]=[\,N\in\mathcal M^{*}\,]$; it returns \SAT{} iff $N_0\in\mathcal M^{*}$, and runs in time
$2^{\mathrm{poly}(n)}$.
\end{theorem}
\begin{proof}
By Lemma~\ref{lem:finiteness} with memoisation the traversal is a depth-first search of a finite
graph, so it terminates. Components are popped in reverse topological order, and by the Tarjan
invariant every child of $C$ outside $C$ has retired when $C$ is popped, one still on the stack
dragging the lowlink of $C$'s root below its index; so $\Out_C$ is final and every pop meets
Lemma~\ref{lem:search}(iii). The deletion loop decreases on a finite set and each step instantiates a
clause of $F$, exact by Lemma~\ref{lem:search}(ii), so it reaches $\nu F_C=\mathcal M^{*}\cap C$, and
by Lemma~\ref{lem:search}(i) pruning preserves every verdict.

For the bound let $m$ be that of Lemma~\ref{lem:finiteness}. As $M$ occupies at most $n$ input
symbols, $\log m=O(n)+O(n^{2})\log M=\mathrm{poly}(n)$. The explored graph has at most $m$ nodes and
$O(nm)$ edges, and memoisation on $\cong_M$-classes expands each once, avoiding the doubly
exponential blow-up of an ancestor-only unfolding. A component takes at most $|C|$ rounds,
each deleting a node or stopping, of one decomposition and $O(n)$ backward reachabilities linear in
the subgraph, and the components partition the nodes, so the run costs $2^{\mathrm{poly}(n)}$.
\end{proof}

\subsection{Soundness and Completeness}\label{ss:sound}

\begin{lemma}\label{lem:zones}
Let $N_0,\dots,N_k$ be a path of $G(\varphi)$ with non-empty zones. \emph{(i)} For every $\nu_k\in
Z_k$ there are $\nu_i\in Z_i$ and delays realising every step, positive at every \ruleref{Delay}
step. \emph{(ii)} If a delay crosses no integer hyperplane $x=b$ with $b\le M$, every subformula of
$\varphi$ has constant truth on the open segment it sweeps. \emph{(iii)} A cycle of $G(\varphi)$ has
a non-Zeno unfolding iff it is time-progressing, and then one whose iterations elapse a constant
positive amount.
\end{lemma}
\begin{proof}
(i) Every zone operation a rule performs (intersection with a single-clock guard, reset,
time-successor, $\cong_M$-canonicalisation) back-propagates, non-empty precisely because the
successor valuation lies in its image~\cite{Bouyer05TimedAutomata}; thread backwards from $\nu_k$,
noting that $\Dopen Z$ and $\Dbd Z$ lie strictly ahead of $Z$ in time. (ii) No clock crosses an
integer at or below $M$, so no integer part changes and no fractional ordering wraps, while clocks
above $M$ are identified by extrapolation; the segment lies in one region, on which TCTL truth is
constant~\cite{AlurCourcoubetisDill1990,Bouyer05TimedAutomata}, and its endpoints are nodes of their
own by the integer split of \ruleref{Delay}. (iii) A cycle elapsing no time takes zero time over
infinitely many traversals; otherwise some clock is neither reset nor extrapolated above $M$, so it
keeps a bound $c\le M$, grows by exactly the delay elapsed, and caps the total at $c$. Conversely (i)
yields a traversal elapsing some $\delta>0$, and the cycle returning to the same node, no rule
distinguishes two visits and no budget depletes, so $\delta$ recurs at every iteration.
\end{proof}

\noindent
\textbf{From a marking to an automaton.} Call a node of $\mathcal M^{*}$ a \emph{state} if it is
saturated, so that its children are the successor steps; the remaining, \emph{expansion} edges
neither elapse time nor change location, and reset a clock only through a \ruleref{launch} or a
\ruleref{Ustoproll}. At most $n+1$ obligations are pending at $N$, the divergence requirement of
\ref{M1} and the until pursuits active there, and \ref{M1}, \ref{M2} give a finite path $\pi_{N,o}$
of $\mathcal M^{*}$ discharging each. A \emph{schedule} selects at every \ORn-node the child
$\pi_{N,o}$ prescribes and moves round-robin to the next pending obligation once $o$ is discharged,
so its memory, the pair (obligation, position on $\pi_{N,o}$), is finite. The witness
$\mathcal A_\varphi$ has as locations the pairs (state, schedule memory), labelled by
$\Gamma_{\mathsf{prop}}$ and carrying $\Inv(N)$, and as edges the scheduled successor steps, guarded
by the zone of the next state reached through expansion edges and resetting the clocks that segment
resets.

\begin{lemma}[Recurrence]\label{lem:recurrence}
Let $\mu=\mathrm{A}(\varphi\,\mathbf{U}_\kI\,\psi)$ with $a=\inf\kI$, $b=\sup\kI<\infty$, and let
$\rho$ be a run of $\mathcal A_\varphi$ through a node carrying $\mu^{\ast}$, launched at $t_0$. Then
\ruleref{Ustoproll} fires on $\rho$ at times $t_0<w_1<w_2<\cdots$ with $w_1-t_0\in\kI$ and
$w_{j+1}-w_j\le b-a$, and $\rho\models\varphi\,\mathbf{U}_\kI\,\psi$ from every position at or after
$t_0$.
\end{lemma}
\begin{proof}
No rule drops $\mu^{\ast}$: \ruleref{Ustoproll} keeps it and \ruleref{Ugo} moves it to $\Ctx$, whence
\ruleref{EX} returns it to $\Gamma$ at every successor. It is carried forever, and with it the
deadline on $x_\mu$, which is $b$ until the first stop and $b-a$ after. Both are at most $M$, so
$x_\mu$ is bounded along $\rho$ and, by Lemma~\ref{lem:search}(ii), every time-progressing cycle
$\rho$ meets resets it. Only \ruleref{Ustoproll} does, \ruleref{Cov} being blocked while the pursuit
is live, and by \ref{M1} the schedule reaches such a cycle and traverses it without end, so the stop
recurs; its first guard is $x_\mu\in\kI$ and each later one is taken under $x_\mu\le b-a$ measured
from the preceding stop, which are the two bounds. Now let $t\ge t_0$ be a position and
$w_j$ the first stop at or after $t+a$. If $j=1$ then $w_1\le t_0+b\le t+b$; otherwise $w_{j-1}<t+a$,
so $w_j\le w_{j-1}+(b-a)<t+b$. Either way $w_j-t\in\kI$, and $\psi$ holds at $w_j$ while
\ruleref{Ugo} asserts $\varphi$ throughout, which is \eqref{eq:sat_u}.
\end{proof}

\begin{theorem}[Soundness]\label{thm:soundness}
If Algorithm~\ref{alg:otf} returns \SAT{} then $\varphi$ is satisfiable, and $\mathcal A_\varphi$ is
a witness for it.
\end{theorem}
\begin{proof}
By Theorem~\ref{thm:alg} the root lies in $\mathcal M^{*}$, so $\mathcal A_\varphi$ is defined, and
Lemma~\ref{lem:zones}(i) supplies valuations and delays along its paths. It is non-blocking: its only
invariants are the deadlines in $\Deadline$, and cap and stop guard derive from the same interval with the
same strictness, so the stop edge is enabled up to the cap and by \ref{M2} the schedule takes it. By
\ref{M1} the schedule reaches a time-progressing cycle, which Lemma~\ref{lem:zones}(iii) unfolds into
a run whose iterations each elapse a fixed $\delta>0$, so $\mathrm{Runs}^\infty$ is non-empty
everywhere, which is all the semantics asks of it.

We show $\mathcal T,(s_0,\mathbf 0)\models\varphi$ for the induced TTS by induction on $\varphi$
along the marking. Propositional goals hold by saturation and closure of contradictory nodes;
$\EX\psi$ obtains a successor satisfying $\psi$ from its \ruleref{EX} child, and $\AX\psi$ holds at
all successors since each child inherits $\Ctx$. For a timed modality demanded at $N$,
\ruleref{Cov} is conjunctive, so $\mathcal M^{*}$ retains both the discharge child, carrying the
remaining goals, and the launch, which establishes $\mu$ with $x_\mu$ reset, so \kI is measured from
this configuration and the pursuit rides every successor, as a goal under $\mathrm{E}$ and as context
under $\mathrm{A}$. An existential until is witnessed by the path \ref{M2} supplies, extended past
its \ruleref{Ustop} by a divergent run; a universal one cannot be evaded, $\Inv(N)$ blocking time at
the cap while the pursuit is imposed on all successors, and when the demand is persistent this is
Lemma~\ref{lem:recurrence}. \ruleref{Rsplit} discharges a release by partitioning the zone
conjunctively: in the window \ruleref{Runf_in} asserts $\psi$ until $\varphi$ releases it, before it
\ruleref{Runf_before} carries the pursuit on, and past it the obligation has lapsed.
Lemma~\ref{lem:zones}(ii) lifts each conclusion to the positions between nodes.
\end{proof}

\begin{theorem}[Completeness]\label{thm:completeness}
If $\varphi$ is non-punctual and satisfiable then Algorithm~\ref{alg:otf} returns \SAT.
\end{theorem}
\begin{proof}
Completeness goes the other way: the realisable nodes of \S\ref{ss:nodes} are themselves a successful
marking. Writing $R$ for them, we show $R\subseteq F(R)$; then $R\subseteq\mathcal M^{*}$ by
Proposition~\ref{prop:gfp}, and $N_0\in R$ as $\varphi$ is satisfiable, so Theorem~\ref{thm:alg}
returns \SAT.

\ref{D1}: a realisable node has a non-empty zone and a consistent $\Gamma_{\mathsf{prop}}$.
\ref{D2}: let $N\in R$ be realised by $\mathcal T$. All children of an \ANDn-rule are realised by the
same $\mathcal T$: \ruleref{Cov} because $\mathcal T$ satisfies $\Gamma\setminus\{\mu\}$ and $\mu$
alike, \ruleref{Rsplit} because its children partition $Z_N$ and realisability is monotone under
narrowing, \ruleref{Delay} and \ruleref{EX} because $\mathcal T$ admits the passage of time and
supplies a successor for every $\EX\psi$, each inheriting $\Ctx$. At an \ORn-rule $\mathcal T$ selects
a child: a disjunct it satisfies for \ruleref{Disj}; for \ruleref{Uunf} and \ruleref{Uunfroll} the
stop if $\psi$ holds now with $x_\mu\in\kI$ and \ruleref{Ugo} otherwise, the witness being later and
$\varphi$ holding now; for \ruleref{Runf_before} and \ruleref{Runf_in} the stop if $\varphi$ has
released the obligation and the recursive child otherwise. The zone operations being exact, no chosen
zone is empty; and that realisability quantifies universally over $Z_N$ is what lets one child answer
a demand re-imposed at a continuum of onsets.
\ref{D3} and \ref{D4}: fix $N\in R$ and a run $\rho$ of $\mathcal T$ from a configuration for $N$. It
is non-Zeno, so no clock stays both unreset and bounded along it, and following $\rho$ through the
choices above gives a path of $R$ from $N$ which by Lemma~\ref{lem:finiteness} repeats a node; that
segment elapses positive time with every clock reset on it or extrapolated above $M$, hence is a
time-progressing cycle of $R$ reached from $N$. An until pursuit active at $N$ is one $\mathcal T$
fulfils at a finite position of $\rho$, which the path following $\rho$ carries it to.

Non-punctuality enters at the repetition step: a node summarises the pending obligations of a
persistent universal until by $x_\mu$ alone, the time since the last witness capped at $b-a$, and by
Lemma~\ref{lem:recurrence} that summary is faithful exactly when witnesses spaced at most $b-a$ apart
answer every onset, which needs $b>a$.
\end{proof}


\begin{figure}[b]
\centering
\begin{tikzpicture}[x=1mm,y=1mm,font=\scriptsize]
\begin{scope}
\draw[->] (14,22) -- (58,22);
\foreach \x/\lab in {20/1,34/2,48/3} \draw (\x,21.4) -- (\x,22.6) node[above=0pt]{\lab};
\node at (13,17) {$0$};   \draw[line width=1.4pt] (20,17) -- (34,17);
\myvdots{13}{14.7}{0.7}{0.1}
\node at (13,11) {$0.4$}; \draw[line width=1.4pt] (25.6,11) -- (39.6,11);
\myvdots{13}{8.7}{0.7}{0.1}
\node at (13,5) {$0.8$};  \draw[line width=1.4pt] (31.2,5) -- (45.2,5);
\myvdots{13}{2.7}{0.7}{0.1}
\draw[dashed] (33,18.5) -- (33,3.5);
\foreach \y in {17,11,5} \fill (33,\y) circle (0.75);
\node at (36,-1) {(a)\ $\kI=[1,2]$};
\end{scope}
\begin{scope}[xshift=52mm]
\draw[->] (14,22) -- (58,22);
\foreach \x/\lab in {20/1,34/2,48/3} \draw (\x,21.4) -- (\x,22.6) node[above=0pt]{\lab};
\node at (13,17) {$0$};   \fill (20,17) circle (0.75);
\myvdots{13}{14.7}{0.7}{0.1}
\node at (13,11) {$0.4$}; \fill (25.6,11) circle (0.75);
\myvdots{13}{8.7}{0.7}{0.1}
\node at (13,5) {$0.8$};  \fill (31.2,5) circle (0.75);
\myvdots{13}{2.7}{0.7}{0.1}
\node at (36,-1) {(b)\ $\kI=[1,1]$};
\end{scope}
\end{tikzpicture}
\caption{Onsets of a persistent demand and their windows: overlapping for a
non-punctual \kI\ \subref*{fig:non_punctual}, pairwise disjoint for a punctual one
\subref*{fig:punctual}.}
\label{fig:punctual_vs_nonpunctual}%
\phantomsubcaption\label{fig:non_punctual}%
\phantomsubcaption\label{fig:punctual}
\end{figure}

\begin{corollary}\label{cor:depthone}
Satisfiability of non-punctual TCTL is \textsc{ExpTime}-complete, and the procedure is sound and
complete for TCTL of temporal nesting depth at most one, punctual intervals included.
\end{corollary}
\begin{proof}
Membership is Theorems~\ref{thm:alg}--\ref{thm:completeness}; hardness is inherited from CTL
satisfiability~\cite{Clarke99} by fixing every interval to $[0,\infty)$. At depth one every argument
is propositional, so no modality is demanded inside the pursuit of another and no fresh onset arises;
the obligations live at a node are exactly the pursuits it carries, recorded in $\Gamma$, $\Ctx$ and
$\Deadline$ with their elapsed time in $Z$, so a repeat agrees on the full obligation state and
Theorem~\ref{thm:completeness} goes through for every \kI.
\end{proof}

\noindent
Fig.~\ref{fig:punctual_vs_nonpunctual} draws the gap: for $\kI=[a,b]$ with $a<b$ the windows of
nearby onsets overlap, so the pending obligations collapse to the earliest deadline and $x_\mu$
records it, whereas for $\kI=[a,a]$ they are disjoint singletons and what a model owes depends on the
exact offsets of every onset still pending.

\noindent
Depth one covers reachability, safety and inevitability, the queries
\textsc{Uppaal}~\cite{BehrmannLarsenYi2004Uppaal} and \textsc{Romeo}~\cite{lime-TACAS-09} are asked
for; the one common query beyond it, leads-to $\mathrm{AG}(p\rightarrow\mathrm{AF}_\kI q)$, stays
covered by Theorem~\ref{thm:completeness} whenever \kI is non-punctual.


%% =====================================================================
%% APPENDIX A -- deferred proofs.  \input this after \bibliography in
%% ETAPS2026.tex, inside an \appendix block:
%%
%%   \appendix
%%   %% =====================================================================
%% DRAFT / extended version of Section 5.  NOT the file \input by
%% ETAPS2026.tex; the shipped version is sec/section5_OTF.tex.
%%
%% This file is the same section with every proof inline and the
%% punctuality figure kept in place (~6 pages).  The shipped version
%% (~4.6 pages) states Proposition 1 and Lemmas 1-4 without proof and
%% points at the appendix; APPENDIX A below is exactly that appendix,
%% ready to \input after the bibliography.
%%
%% Design notes are at the end of the file.
%% =====================================================================

\section{The Decision Procedure}\label{sec:correct}

Section~\ref{sec:tableau} defines success by a condition on a marking of a graph that no
implementation can afford to build. We make that graph finite and identify the largest successful
marking with a greatest fixpoint (\S\ref{ss:gfp}), give a procedure deciding membership of the root
in it while building only what the search demands (\S\ref{ss:otf}), and prove that procedure sound
for TCTL and complete for the non-punctual fragment (\S\ref{ss:sound}). Throughout $\varphi$ is the
input, $n=|\FL(\varphi)|$, and $M$ the largest constant in its intervals. 

\subsection{The tableau graph and its greatest marking}\label{ss:gfp}

The calculus leaves the order of rule application open; the procedure fixes one. Let
$\FL^{\bullet}(\varphi)$ extend $\FL(\varphi)$ with the pursuit markers $\mu^\bullet,\mu^{\ast}$ and
with $\EX\chi,\AX\chi$ for $\chi$ in it, so that it holds every goal the rules produce and has size
$O(n)$, and fix a linear order on it. Then $\Expand{N}$ is empty if $N$ is closed, and otherwise
fires on the least goal of $\Gamma$ it fits the first applicable rule in the order \ruleref{Conj},
\ruleref{Disj}, \ruleref{AX}, \ruleref{Roll}, \ruleref{Cov}, \ruleref{Rsplit},
\ruleref{Runf_before}, \ruleref{Runf_in}, \ruleref{Uunf}, \ruleref{Uunfroll}; when none applies, $N$
is saturated and $\Expand{N}$ returns the \ruleref{Delay} and \ruleref{EX} children together, both
being conjunctive. Two conventions make this a function of the node: the clock a \ruleref{launch}
allocates is determined by $\mu$ rather than drawn fresh, sound because \ruleref{Cov} blocks a second
launch while a pursuit of $\mu$ is live, so $\clockSet=\{x_\mu\mid\mu$ timed in $\FL(\varphi)\}$ has
size $O(n)$; and acceptance reads clock resets and stop branches off the graph, both properties of an
edge, so each node records the clocks its rule reset and the modality it discharged.

\begin{definition}[Tableau graph]\label{def:graph}
$G(\varphi)$ is the least set of nodes containing $N_0=\tup{\{\varphi\}}{\emptyset}{\{\mathbf
0\}}{\emptyset}{\emptyset}$ and closed under $\Expand{}$, two nodes being identified when they agree
componentwise and their zones agree modulo $\cong_M$, with an edge $N\to N'$ for each
$N'\in\Expand{N}$.
\end{definition}

\begin{lemma}[Finiteness]\label{lem:finiteness}
$|G(\varphi)|\le 2^{O(n)}\cdot M^{O(n^{2})}$, and every node has $O(n)$ children.
\end{lemma}
\begin{proof}
$\Gamma$, $\Ctx$, the reset set and the discharge marker range over subsets of
$\FL^{\bullet}(\varphi)$ or of $\clockSet$; $\obligations$ is fixed by its domain; and $\Deadline$ gives each
$x_\mu$ one of the bounds $\sup\kI$, $\sup\kI-\inf\kI$ with a strictness flag. Each contributes
$2^{O(n)}$. A canonical $\cong_M$-zone over $K\le n+1$ clocks is a $(K{+}1)^{2}$ DBM with entries
$\infty$ or in $\{-M,\dots,M\}\times\{<,\le\}$, giving $(4M{+}3)^{(n+2)^{2}}=M^{O(n^{2})}$ zones.
Every rule has at most three children except \ruleref{EX}, with one per $\EX\psi\in\Gamma$.
\end{proof}

\noindent
Identification up to $\cong_M$ dispenses with a termination argument for expansion, and so lets the
calculus afford rules that re-impose their own premise: a chain of applications that would regress
forever revisits a node and closes a cycle, which crosses no \ruleref{Delay} step and is discarded by
\ref{M1}.

A marking contains no closed node, at least one child of each of its \ORn-nodes and all children of
each of its \ANDn-nodes, and is successful when every member satisfies \ref{M1} and \ref{M2}. Both
clauses quantify over paths \emph{inside} the marking, so a successful marking is a set that survives
the deletions it prescribes. For $A\subseteq G(\varphi)$ let $F(A)$ be $A$ minus:
\begin{enumerate}[start=1,label={D\arabic*)}, itemindent=1.6em, ref=(D\arabic*), itemsep=0.02em]
\item\label{D1} the closed nodes;
\item\label{D2} the \ORn-nodes with no child in $A$, and the \ANDn-nodes with a child outside $A$;
\item\label{D3} the nodes with no $A$-path to a node on a time-progressing cycle of $A$;
\item\label{D4} for each until $\mu$, the nodes pursuing $\mu$ from which no $A$-path carrying the
      pursuit throughout reaches a stop for $\mu$ taken where the zone meets $x_\mu\in\kI$.
\end{enumerate}

\begin{proposition}\label{prop:gfp}
$F$ is monotone, and $A$ is a successful marking iff $A\subseteq F(A)$. Hence the successful markings
have a greatest element $\mathcal M^{*}=\nu F$, and the tableau succeeds on $\varphi$ iff
$N_0\in\mathcal M^{*}$.
\end{proposition}
\begin{proof}
\ref{D1} ignores $A$; a node kept by \ref{D2} has its witnessing child, or all its children, in $A$
and hence in any superset; \ref{D3} and \ref{D4} ask for a path in $A$, which survives enlargement.
So $F$ is monotone on the finite lattice $2^{G(\varphi)}$, and $A\subseteq F(A)$ unfolds to the
clauses defining a successful marking. Knaster--Tarski gives a greatest post-fixpoint, equal to $\nu
F$ and to the union of all post-fixpoints, and the tableau succeeds iff one of them contains $N_0$.
\end{proof}

\subsection{The on-the-fly procedure}\label{ss:otf}

Algorithm~\ref{alg:otf} decides $N_0\in\mathcal M^{*}$ without ever holding $G(\varphi)$.
$\Expand{}$ applies the rules the first time the search reaches a node and returns the recorded
children of a $\cong_M$-equal node otherwise, so the graph exists only as far as the search has
demanded it. Tarjan's indices, used as a black box, ride along; a back edge yields a lowlink update
and nothing else, a node still on the stack having no verdict and none being guessed for it. The
third line of the loop is where the search stops at the first witness: a child settles its parent
when the parent is an \ORn-node and the child holds, or an \ANDn-node and it fails, which is what
$\val[w]=[\,v$ is $\ORn\,]$ says, and the parent retires with its remaining children unbuilt; the guard $\mathrm{top}(\Stack)=v$ makes this well defined, an earlier child
possibly having left members of $v$'s component above $v$ on the stack. A component is decided the
moment it closes, against the verdicts of the children it leaves: writing
$\Out_C=\{w\notin C\mid w\in\ch[v]$ for some $v\in C$, $\val[w]=\top\}$ for those already accepted,
$\nu F_C$ is the largest $A\subseteq C$ with $A\subseteq F(A\cup\Out_C)$, computed by deletion from
$A=C$, each round applying \ref{D1} and \ref{D2} pointwise and then the backward reachabilities
\ref{D3}, \ref{D4} over the survivors. A deletion can disconnect $A$, so the decomposition is
recomputed each round.

\begin{algorithm}[!tb]
\footnotesize
\DontPrintSemicolon
\SetAlgoLined
\SetKwProg{Fn}{procedure}{}{end}

\Fn{\Decide{$\varphi$}}{
  $\val[v]\gets{?}$ and $\idx[v]\gets\bot$ for all $v$; \ $\Stack\gets\varepsilon$; \ $\cnt\gets0$\;
  \Explore{$N_0$} \ \textbf{where} $N_0=\tup{\{\varphi\}}{\emptyset}{\{\mathbf
    0\}}{\emptyset}{\emptyset}$; \quad \Return{\SAT{} \textnormal{iff} $\val[N_0]=\top$}\;
}
\BlankLine
\Fn{\Explore{$v$}}{
  $\idx[v]\gets\low[v]\gets\cnt{+}{+}$; \ $\mathrm{push}(\Stack,v)$; \
  $\ch[v]\gets\Expand{v}$ \tcp*{memoised on $\cong_M$}
  \ForEach{$w\in\ch[v]$}{
    \lIf{$\idx[w]=\bot$}{\Explore{$w$}; \ $\low[v]\gets\min(\low[v],\low[w])$}
    \lElseIf{$w\in\Stack$}{$\low[v]\gets\min(\low[v],\idx[w])$}
    \lIf{$\val[w]=[\,v\text{ is }\ORn\,]$ \textnormal{\textbf{and}} $\mathrm{top}(\Stack)=v$}
        {$\val[v]\gets\val[w]$; \ $\mathrm{pop}(\Stack)$; \Return}
  }
  \lIf{$\low[v]=\idx[v]$}{$C\gets$ pop $\Stack$ down to $v$; \
    $\val[u]\gets[\,u\in\nu F_C\,]$ for $u\in C$}
}
\caption{On-the-fly satisfiability.}\label{alg:otf}
\end{algorithm}

\begin{lemma}\label{lem:search}
\emph{(i)} If $\val[w]$ is final and settles $v$, then $v\in\mathcal M^{*}$ iff $w\in\mathcal M^{*}$,
and leaving the remaining children of $v$ unexpanded changes $\mathcal M^{*}$ at no node.
\emph{(ii)} A strongly connected $T$ with an internal edge carries a time-progressing cycle iff some
node of $T$ is a \ruleref{Delay}-successor and every clock is reset somewhere in $T$ or unbounded
everywhere in $T$.
\emph{(iii)} If $C$ is a maximal strongly connected set with $\val[c]=[\,c\in\mathcal M^{*}\,]$ for
every child $c\notin C$ of a node of $C$, then $\mathcal M^{*}\cap C=\nu F_C$.
\end{lemma}
\begin{proof}
(i) The first claim is \ref{D2}. The explored subgraph is closed under $\Expand{}$ except at retired
nodes, so a path of $\mathcal M^{*}$ using an unexpanded child of $v$ passes through $v$; if
$v\notin\mathcal M^{*}$ there is no such path, and if $v\in\mathcal M^{*}$ then $v$ satisfies
\ref{M1}, and \ref{M2} for each pursuit active at it, so a node reaching $v$ inherits both witnesses
by prefixing the path to $v$, which carries a pursuit exactly when it is active at $v$.

(ii) $(\Leftarrow)$ Concatenating paths in $T$ gives a cycle through the \ruleref{Delay}-successor,
which therefore elapses positive time, and through one reset node of every clock $T$ resets; the rest
are unbounded throughout. $(\Rightarrow)$ A time-progressing cycle in $T$ elapses positive time,
hence crosses a \ruleref{Delay} step; and a clock reset nowhere in $T$ is forced unbounded on that
cycle, whence unbounded on all of $T$, since a path inside $T$ never resets it and its value, already
above $M$, only grows.

(iii) $C$ being maximal, a path of $\mathcal M^{*}$ from $N\in C$ either stays in $C$, and a witness
for \ref{M1} or \ref{M2} then lies inside $C$; or it leaves $C$ once and for all at some
$c\in\mathcal M^{*}\setminus C$, which by hypothesis lies in $\Out_C$ and itself satisfies \ref{M1},
and \ref{M2} for every pursuit active at it, so prefixing the path from $N$ to $c$ turns either
witness into one for $N$. Hence $\mathcal M^{*}\cap C\subseteq F(\mathcal M^{*}\cap C\cup\Out_C)$.
Conversely any $A\subseteq C$ with $A\subseteq F(A\cup\Out_C)$ extends by the nodes of
$\mathcal M^{*}$ reachable from $\Out_C$ to a post-fixpoint of $F$, which Proposition~\ref{prop:gfp}
places inside $\mathcal M^{*}$. The greatest elements coincide.
\end{proof}

\begin{theorem}[Adequacy and complexity]\label{thm:alg}
Algorithm~\ref{alg:otf} terminates and retires every node it builds with
$\val[N]=[\,N\in\mathcal M^{*}\,]$; it returns \SAT{} iff $N_0\in\mathcal M^{*}$, and runs in time
$2^{\mathrm{poly}(n)}$.
\end{theorem}
\begin{proof}
By Lemma~\ref{lem:finiteness} with memoisation the traversal is a depth-first search of a finite
graph, so it terminates. Components are popped in reverse topological order, and by the Tarjan
invariant every child of $C$ outside $C$ has retired when $C$ is popped, one still on the stack
dragging the lowlink of $C$'s root below its index; so $\Out_C$ is final and every pop meets
Lemma~\ref{lem:search}(iii). The deletion loop decreases on a finite set and each step instantiates a
clause of $F$, exact by Lemma~\ref{lem:search}(ii), so it reaches $\nu F_C=\mathcal M^{*}\cap C$, and
by Lemma~\ref{lem:search}(i) pruning preserves every verdict.

For the bound let $m$ be that of Lemma~\ref{lem:finiteness}. As $M$ occupies at most $n$ input
symbols, $\log m=O(n)+O(n^{2})\log M=\mathrm{poly}(n)$. The explored graph has at most $m$ nodes and
$O(nm)$ edges, and memoisation on $\cong_M$-classes expands each once, avoiding the doubly
exponential blow-up of an ancestor-only unfolding. A component takes at most $|C|$ rounds,
each deleting a node or stopping, of one decomposition and $O(n)$ backward reachabilities linear in
the subgraph, and the components partition the nodes, so the run costs $2^{\mathrm{poly}(n)}$.
\end{proof}

\subsection{Soundness and Completeness}\label{ss:sound}

\begin{lemma}\label{lem:zones}
Let $N_0,\dots,N_k$ be a path of $G(\varphi)$ with non-empty zones. \emph{(i)} For every $\nu_k\in
Z_k$ there are $\nu_i\in Z_i$ and delays realising every step, positive at every \ruleref{Delay}
step. \emph{(ii)} If a delay crosses no integer hyperplane $x=b$ with $b\le M$, every subformula of
$\varphi$ has constant truth on the open segment it sweeps. \emph{(iii)} A cycle of $G(\varphi)$ has
a non-Zeno unfolding iff it is time-progressing, and then one whose iterations elapse a constant
positive amount.
\end{lemma}
\begin{proof}
(i) Every zone operation a rule performs (intersection with a single-clock guard, reset,
time-successor, $\cong_M$-canonicalisation) back-propagates, non-empty precisely because the
successor valuation lies in its image~\cite{Bouyer05TimedAutomata}; thread backwards from $\nu_k$,
noting that $\Dopen Z$ and $\Dbd Z$ lie strictly ahead of $Z$ in time. (ii) No clock crosses an
integer at or below $M$, so no integer part changes and no fractional ordering wraps, while clocks
above $M$ are identified by extrapolation; the segment lies in one region, on which TCTL truth is
constant~\cite{AlurCourcoubetisDill1990,Bouyer05TimedAutomata}, and its endpoints are nodes of their
own by the integer split of \ruleref{Delay}. (iii) A cycle elapsing no time takes zero time over
infinitely many traversals; otherwise some clock is neither reset nor extrapolated above $M$, so it
keeps a bound $c\le M$, grows by exactly the delay elapsed, and caps the total at $c$. Conversely (i)
yields a traversal elapsing some $\delta>0$, and the cycle returning to the same node, no rule
distinguishes two visits and no budget depletes, so $\delta$ recurs at every iteration.
\end{proof}

\noindent
\textbf{From a marking to an automaton.} Call a node of $\mathcal M^{*}$ a \emph{state} if it is
saturated, so that its children are the successor steps; the remaining, \emph{expansion} edges
neither elapse time nor change location, and reset a clock only through a \ruleref{launch} or a
\ruleref{Ustoproll}. At most $n+1$ obligations are pending at $N$, the divergence requirement of
\ref{M1} and the until pursuits active there, and \ref{M1}, \ref{M2} give a finite path $\pi_{N,o}$
of $\mathcal M^{*}$ discharging each. A \emph{schedule} selects at every \ORn-node the child
$\pi_{N,o}$ prescribes and moves round-robin to the next pending obligation once $o$ is discharged,
so its memory, the pair (obligation, position on $\pi_{N,o}$), is finite. The witness
$\mathcal A_\varphi$ has as locations the pairs (state, schedule memory), labelled by
$\Gamma_{\mathsf{prop}}$ and carrying $\Inv(N)$, and as edges the scheduled successor steps, guarded
by the zone of the next state reached through expansion edges and resetting the clocks that segment
resets.

\begin{lemma}[Recurrence]\label{lem:recurrence}
Let $\mu=\mathrm{A}(\varphi\,\mathbf{U}_\kI\,\psi)$ with $a=\inf\kI$, $b=\sup\kI<\infty$, and let
$\rho$ be a run of $\mathcal A_\varphi$ through a node carrying $\mu^{\ast}$, launched at $t_0$. Then
\ruleref{Ustoproll} fires on $\rho$ at times $t_0<w_1<w_2<\cdots$ with $w_1-t_0\in\kI$ and
$w_{j+1}-w_j\le b-a$, and $\rho\models\varphi\,\mathbf{U}_\kI\,\psi$ from every position at or after
$t_0$.
\end{lemma}
\begin{proof}
No rule drops $\mu^{\ast}$: \ruleref{Ustoproll} keeps it and \ruleref{Ugo} moves it to $\Ctx$, whence
\ruleref{EX} returns it to $\Gamma$ at every successor. It is carried forever, and with it the
deadline on $x_\mu$, which is $b$ until the first stop and $b-a$ after. Both are at most $M$, so
$x_\mu$ is bounded along $\rho$ and, by Lemma~\ref{lem:search}(ii), every time-progressing cycle
$\rho$ meets resets it. Only \ruleref{Ustoproll} does, \ruleref{Cov} being blocked while the pursuit
is live, and by \ref{M1} the schedule reaches such a cycle and traverses it without end, so the stop
recurs; its first guard is $x_\mu\in\kI$ and each later one is taken under $x_\mu\le b-a$ measured
from the preceding stop, which are the two bounds. Now let $t\ge t_0$ be a position and
$w_j$ the first stop at or after $t+a$. If $j=1$ then $w_1\le t_0+b\le t+b$; otherwise $w_{j-1}<t+a$,
so $w_j\le w_{j-1}+(b-a)<t+b$. Either way $w_j-t\in\kI$, and $\psi$ holds at $w_j$ while
\ruleref{Ugo} asserts $\varphi$ throughout, which is \eqref{eq:sat_u}.
\end{proof}

\begin{theorem}[Soundness]\label{thm:soundness}
If Algorithm~\ref{alg:otf} returns \SAT{} then $\varphi$ is satisfiable, and $\mathcal A_\varphi$ is
a witness for it.
\end{theorem}
\begin{proof}
By Theorem~\ref{thm:alg} the root lies in $\mathcal M^{*}$, so $\mathcal A_\varphi$ is defined, and
Lemma~\ref{lem:zones}(i) supplies valuations and delays along its paths. It is non-blocking: its only
invariants are the deadlines in $\Deadline$, and cap and stop guard derive from the same interval with the
same strictness, so the stop edge is enabled up to the cap and by \ref{M2} the schedule takes it. By
\ref{M1} the schedule reaches a time-progressing cycle, which Lemma~\ref{lem:zones}(iii) unfolds into
a run whose iterations each elapse a fixed $\delta>0$, so $\mathrm{Runs}^\infty$ is non-empty
everywhere, which is all the semantics asks of it.

We show $\mathcal T,(s_0,\mathbf 0)\models\varphi$ for the induced TTS by induction on $\varphi$
along the marking. Propositional goals hold by saturation and closure of contradictory nodes;
$\EX\psi$ obtains a successor satisfying $\psi$ from its \ruleref{EX} child, and $\AX\psi$ holds at
all successors since each child inherits $\Ctx$. For a timed modality demanded at $N$,
\ruleref{Cov} is conjunctive, so $\mathcal M^{*}$ retains both the discharge child, carrying the
remaining goals, and the launch, which establishes $\mu$ with $x_\mu$ reset, so \kI is measured from
this configuration and the pursuit rides every successor, as a goal under $\mathrm{E}$ and as context
under $\mathrm{A}$. An existential until is witnessed by the path \ref{M2} supplies, extended past
its \ruleref{Ustop} by a divergent run; a universal one cannot be evaded, $\Inv(N)$ blocking time at
the cap while the pursuit is imposed on all successors, and when the demand is persistent this is
Lemma~\ref{lem:recurrence}. \ruleref{Rsplit} discharges a release by partitioning the zone
conjunctively: in the window \ruleref{Runf_in} asserts $\psi$ until $\varphi$ releases it, before it
\ruleref{Runf_before} carries the pursuit on, and past it the obligation has lapsed.
Lemma~\ref{lem:zones}(ii) lifts each conclusion to the positions between nodes.
\end{proof}

\begin{theorem}[Completeness]\label{thm:completeness}
If $\varphi$ is non-punctual and satisfiable then Algorithm~\ref{alg:otf} returns \SAT.
\end{theorem}
\begin{proof}
Completeness goes the other way: the realisable nodes of \S\ref{ss:nodes} are themselves a successful
marking. Writing $R$ for them, we show $R\subseteq F(R)$; then $R\subseteq\mathcal M^{*}$ by
Proposition~\ref{prop:gfp}, and $N_0\in R$ as $\varphi$ is satisfiable, so Theorem~\ref{thm:alg}
returns \SAT.

\ref{D1}: a realisable node has a non-empty zone and a consistent $\Gamma_{\mathsf{prop}}$.
\ref{D2}: let $N\in R$ be realised by $\mathcal T$. All children of an \ANDn-rule are realised by the
same $\mathcal T$: \ruleref{Cov} because $\mathcal T$ satisfies $\Gamma\setminus\{\mu\}$ and $\mu$
alike, \ruleref{Rsplit} because its children partition $Z_N$ and realisability is monotone under
narrowing, \ruleref{Delay} and \ruleref{EX} because $\mathcal T$ admits the passage of time and
supplies a successor for every $\EX\psi$, each inheriting $\Ctx$. At an \ORn-rule $\mathcal T$ selects
a child: a disjunct it satisfies for \ruleref{Disj}; for \ruleref{Uunf} and \ruleref{Uunfroll} the
stop if $\psi$ holds now with $x_\mu\in\kI$ and \ruleref{Ugo} otherwise, the witness being later and
$\varphi$ holding now; for \ruleref{Runf_before} and \ruleref{Runf_in} the stop if $\varphi$ has
released the obligation and the recursive child otherwise. The zone operations being exact, no chosen
zone is empty; and that realisability quantifies universally over $Z_N$ is what lets one child answer
a demand re-imposed at a continuum of onsets.
\ref{D3} and \ref{D4}: fix $N\in R$ and a run $\rho$ of $\mathcal T$ from a configuration for $N$. It
is non-Zeno, so no clock stays both unreset and bounded along it, and following $\rho$ through the
choices above gives a path of $R$ from $N$ which by Lemma~\ref{lem:finiteness} repeats a node; that
segment elapses positive time with every clock reset on it or extrapolated above $M$, hence is a
time-progressing cycle of $R$ reached from $N$. An until pursuit active at $N$ is one $\mathcal T$
fulfils at a finite position of $\rho$, which the path following $\rho$ carries it to.

Non-punctuality enters at the repetition step: a node summarises the pending obligations of a
persistent universal until by $x_\mu$ alone, the time since the last witness capped at $b-a$, and by
Lemma~\ref{lem:recurrence} that summary is faithful exactly when witnesses spaced at most $b-a$ apart
answer every onset, which needs $b>a$.
\end{proof}


\begin{figure}[b]
\centering
\begin{tikzpicture}[x=1mm,y=1mm,font=\scriptsize]
\begin{scope}
\draw[->] (14,22) -- (58,22);
\foreach \x/\lab in {20/1,34/2,48/3} \draw (\x,21.4) -- (\x,22.6) node[above=0pt]{\lab};
\node at (13,17) {$0$};   \draw[line width=1.4pt] (20,17) -- (34,17);
\myvdots{13}{14.7}{0.7}{0.1}
\node at (13,11) {$0.4$}; \draw[line width=1.4pt] (25.6,11) -- (39.6,11);
\myvdots{13}{8.7}{0.7}{0.1}
\node at (13,5) {$0.8$};  \draw[line width=1.4pt] (31.2,5) -- (45.2,5);
\myvdots{13}{2.7}{0.7}{0.1}
\draw[dashed] (33,18.5) -- (33,3.5);
\foreach \y in {17,11,5} \fill (33,\y) circle (0.75);
\node at (36,-1) {(a)\ $\kI=[1,2]$};
\end{scope}
\begin{scope}[xshift=52mm]
\draw[->] (14,22) -- (58,22);
\foreach \x/\lab in {20/1,34/2,48/3} \draw (\x,21.4) -- (\x,22.6) node[above=0pt]{\lab};
\node at (13,17) {$0$};   \fill (20,17) circle (0.75);
\myvdots{13}{14.7}{0.7}{0.1}
\node at (13,11) {$0.4$}; \fill (25.6,11) circle (0.75);
\myvdots{13}{8.7}{0.7}{0.1}
\node at (13,5) {$0.8$};  \fill (31.2,5) circle (0.75);
\myvdots{13}{2.7}{0.7}{0.1}
\node at (36,-1) {(b)\ $\kI=[1,1]$};
\end{scope}
\end{tikzpicture}
\caption{Onsets of a persistent demand and their windows: overlapping for a
non-punctual \kI\ \subref*{fig:non_punctual}, pairwise disjoint for a punctual one
\subref*{fig:punctual}.}
\label{fig:punctual_vs_nonpunctual}%
\phantomsubcaption\label{fig:non_punctual}%
\phantomsubcaption\label{fig:punctual}
\end{figure}

\begin{corollary}\label{cor:depthone}
Satisfiability of non-punctual TCTL is \textsc{ExpTime}-complete, and the procedure is sound and
complete for TCTL of temporal nesting depth at most one, punctual intervals included.
\end{corollary}
\begin{proof}
Membership is Theorems~\ref{thm:alg}--\ref{thm:completeness}; hardness is inherited from CTL
satisfiability~\cite{Clarke99} by fixing every interval to $[0,\infty)$. At depth one every argument
is propositional, so no modality is demanded inside the pursuit of another and no fresh onset arises;
the obligations live at a node are exactly the pursuits it carries, recorded in $\Gamma$, $\Ctx$ and
$\Deadline$ with their elapsed time in $Z$, so a repeat agrees on the full obligation state and
Theorem~\ref{thm:completeness} goes through for every \kI.
\end{proof}

\noindent
Fig.~\ref{fig:punctual_vs_nonpunctual} draws the gap: for $\kI=[a,b]$ with $a<b$ the windows of
nearby onsets overlap, so the pending obligations collapse to the earliest deadline and $x_\mu$
records it, whereas for $\kI=[a,a]$ they are disjoint singletons and what a model owes depends on the
exact offsets of every onset still pending.

\noindent
Depth one covers reachability, safety and inevitability, the queries
\textsc{Uppaal}~\cite{BehrmannLarsenYi2004Uppaal} and \textsc{Romeo}~\cite{lime-TACAS-09} are asked
for; the one common query beyond it, leads-to $\mathrm{AG}(p\rightarrow\mathrm{AF}_\kI q)$, stays
covered by Theorem~\ref{thm:completeness} whenever \kI is non-punctual.


%% =====================================================================
%% APPENDIX A -- deferred proofs.  \input this after \bibliography in
%% ETAPS2026.tex, inside an \appendix block:
%%
%%   \appendix
%%   \input{sec/section5_OTF_draft.tex}   % only the part below
%% =====================================================================
\iffalse
\section{Deferred proofs}\label{app:proofs}

\begin{proof}[of Lemma~\ref{lem:finiteness}]
$\Gamma$, $\Ctx$, the reset set and the discharge marker range over subsets of
$\FL^{\bullet}(\varphi)$ or of $\clockSet$; $\obligations$ is fixed by its domain; and $\Deadline$ gives each
$x_\mu$ one of the bounds $\sup\kI$, $\sup\kI-\inf\kI$ with a strictness flag. Each contributes
$2^{O(n)}$. A canonical $\cong_M$-zone over $K\le n+1$ clocks is a $(K{+}1)^{2}$ DBM with entries
$\infty$ or in $\{-M,\dots,M\}\times\{<,\le\}$, giving $(4M{+}3)^{(n+2)^{2}}=M^{O(n^{2})}$ zones.
Every rule has at most three children except \ruleref{EX}, with one per $\EX\psi\in\Gamma$.
\end{proof}

\begin{proof}[of Proposition~\ref{prop:gfp}]
\ref{D1} ignores $A$; a node kept by \ref{D2} has its witnessing child, or all its children, in $A$
and hence in any superset; \ref{D3} and \ref{D4} ask for a path in $A$, which survives enlargement.
So $F$ is monotone on the finite lattice $2^{G(\varphi)}$, and $A\subseteq F(A)$ unfolds to the
clauses defining a successful marking. Knaster--Tarski gives a greatest post-fixpoint, equal to $\nu
F$ and to the union of all post-fixpoints, and the tableau succeeds iff one of them contains $N_0$.
\end{proof}

\begin{proof}[of Lemma~\ref{lem:search}]
(i) The first claim is \ref{D2}. The explored subgraph is closed under $\Expand{}$ except at retired
nodes, so a path of $\mathcal M^{*}$ using an unexpanded child of $v$ passes through $v$; if
$v\notin\mathcal M^{*}$ there is no such path, and if $v\in\mathcal M^{*}$ then $v$ satisfies
\ref{M1}, and \ref{M2} for each pursuit active at it, so a node reaching $v$ inherits both witnesses
by prefixing the path to $v$, which carries a pursuit exactly when it is active at $v$.

(ii) $(\Leftarrow)$ Concatenating paths in $T$ gives a cycle through the \ruleref{Delay}-successor,
which therefore elapses positive time, and through one reset node of every clock $T$ resets; the rest
are unbounded throughout. $(\Rightarrow)$ A time-progressing cycle in $T$ elapses positive time,
hence crosses a \ruleref{Delay} step; and a clock reset nowhere in $T$ is forced unbounded on that
cycle, whence unbounded on all of $T$, since a path inside $T$ never resets it and its value, already
above $M$, only grows.

(iii) $C$ being maximal, a path of $\mathcal M^{*}$ from $N\in C$ either stays in $C$, and a witness
for \ref{M1} or \ref{M2} then lies inside $C$; or it leaves $C$ once and for all at some
$c\in\mathcal M^{*}\setminus C$, which by hypothesis lies in $\Out_C$ and itself satisfies \ref{M1},
and \ref{M2} for every pursuit active at it, so prefixing the path from $N$ to $c$ turns either
witness into one for $N$. Hence $\mathcal M^{*}\cap C\subseteq F(\mathcal M^{*}\cap C\cup\Out_C)$.
Conversely any $A\subseteq C$ with $A\subseteq F(A\cup\Out_C)$ extends by the nodes of
$\mathcal M^{*}$ reachable from $\Out_C$ to a post-fixpoint of $F$, which Proposition~\ref{prop:gfp}
places inside $\mathcal M^{*}$. The greatest elements coincide.
\end{proof}

\begin{proof}[of Lemma~\ref{lem:zones}]
(i) Every zone operation a rule performs (intersection with a single-clock guard, reset,
time-successor, $\cong_M$-canonicalisation) back-propagates, non-empty precisely because the
successor valuation lies in its image~\cite{Bouyer05TimedAutomata}; thread backwards from $\nu_k$,
noting that $\Dopen Z$ and $\Dbd Z$ lie strictly ahead of $Z$ in time. (ii) No clock crosses an
integer at or below $M$, so no integer part changes and no fractional ordering wraps, while clocks
above $M$ are identified by extrapolation; the segment lies in one region, on which TCTL truth is
constant~\cite{AlurCourcoubetisDill1990,Bouyer05TimedAutomata}, and its endpoints are nodes of their
own by the integer split of \ruleref{Delay}. (iii) A cycle elapsing no time takes zero time over
infinitely many traversals; otherwise some clock is neither reset nor extrapolated above $M$, so it
keeps a bound $c\le M$, grows by exactly the delay elapsed, and caps the total at $c$. Conversely (i)
yields a traversal elapsing some $\delta>0$, and the cycle returning to the same node, no rule
distinguishes two visits and no budget depletes, so $\delta$ recurs at every iteration.
\end{proof}

\begin{proof}[of Lemma~\ref{lem:recurrence}]
No rule drops $\mu^{\ast}$: \ruleref{Ustoproll} keeps it and \ruleref{Ugo} moves it to $\Ctx$, whence
\ruleref{EX} returns it to $\Gamma$ at every successor. It is carried forever, and with it the
deadline on $x_\mu$, which is $b$ until the first stop and $b-a$ after. Both are at most $M$, so
$x_\mu$ is bounded along $\rho$ and, by Lemma~\ref{lem:search}(ii), every time-progressing cycle
$\rho$ meets resets it. Only \ruleref{Ustoproll} does, \ruleref{Cov} being blocked while the pursuit
is live, and by \ref{M1} the schedule reaches such a cycle and traverses it without end, so the stop
recurs; its first guard is $x_\mu\in\kI$ and each later one is taken under $x_\mu\le b-a$ measured
from the preceding stop, which are the two bounds. Now let $t\ge t_0$ be a position and
$w_j$ the first stop at or after $t+a$. If $j=1$ then $w_1\le t_0+b\le t+b$; otherwise $w_{j-1}<t+a$,
so $w_j\le w_{j-1}+(b-a)<t+b$. Either way $w_j-t\in\kI$, and $\psi$ holds at $w_j$ while
\ruleref{Ugo} asserts $\varphi$ throughout, which is \eqref{eq:sat_u}.
\end{proof}
\fi


%% =====================================================================
%% DESIGN NOTES (not for the paper)
%% =====================================================================
%
% 1. WHY THE OLD SECTION COULD NOT BE EDITED IN PLACE.
%    sec/section5_OTF_old.tex is written against an earlier calculus: it
%    cites \ruleref{Covstar} and the satisfaction sets $\Sat_\Ctx(\mu)$.
%    Neither exists in sec/section4_nosat.tex, whose cover rule is the
%    conjunctive discharge/launch pair.  Lemma lem:sat of the old file,
%    its longest, related $\Sat_\Ctx$ to truth and has no referent any
%    more; its work is now done by the OR/AND case analysis inside
%    Theorem thm:completeness and by Lemma lem:recurrence.  The
%    $\Sat_\Ctx$ paragraph of the old complexity proof is likewise gone.
%
% 2. WHAT IS NEW RELATIVE TO THE OLD SECTION.
%    - Definition def:graph and the expansion order.  The old text never
%      said what the children of a node are; without a deterministic
%      \Expand{} the algorithm is not well defined and lem:finiteness
%      counts nothing.
%    - The observation that no local-termination lemma is needed:
%      identifying nodes up to \cong_M turns a regressing rule chain into
%      a cycle, and a cycle with no (Delay) step fails (M1).  This is why
%      the calculus can afford rules that re-impose their own premise
%      ((rec) of (U*), and (R^split)).
%    - Proposition prop:gfp, stated and proved rather than asserted.
%    - Lemma lem:recurrence.  This argument existed only as a
%      commented-out paragraph in section4_nosat.tex (lines 489-495).  It
%      is the semantic justification of the b-a retuning in (rec).
%      Without it, soundness for a persistent universal until is not
%      proved.
%    - The schedule / finite memory in the model extraction.  The old
%      proof said "its locations are the nodes of M*", which is wrong for
%      an OR-node with two surviving children: the automaton would
%      contain both, and a universal obligation could fail along the one
%      the marking did not use to establish the disjunction.  The
%      round-robin schedule is the standard fulfilling-path fix.
%    - Completeness restructured as R \subseteq F(R) for R the realisable
%      nodes.  Shorter, and it finally uses the realisability definition
%      of \S4.1, which the old proof did not.
%
% 3. FAITHFULNESS TO Implementation/src/tableau.cpp.
%    - Explore matches Engine::decide_scc (:1543); the settles test with
%      the top-of-stack guard is settles (:1571).
%    - \nu F_C matches evaluate_scc (:1284): optimistic start, inner
%      local-test fixpoint, then query_progress (:1392) and query_until
%      (:1435), outer loop until no deletion.
%    - Lemma lem:search(ii) is progressing() (:1368): union of reset
%      clocks, intersection of unbounded clocks, at least one (Delay).
%    - Back edges get a lowlink update and no verdict (:1584).  The notes
%      record this as load-bearing: optimistic back edges produced 82
%      soundness violations.
%    - "only the clocks of live pursuits and deadlines" is
%      collapse_dead_clocks (:698), an exact active-clock reduction.
%      (That sentence was cut from the shipped version for space.)
%    - The reset/discharge markers in node identity are just_reset and
%      discharges_until in acc_key (:1149-1157).
%    NOT claimed here, deliberately: the prototype still computes
%    satisfaction sets S_X for its cover rule and never populates
%    Node::X.  Both are artefacts of the older calculus.  Align the
%    prototype with \S4 before the evaluation numbers are cited.
%
% 4. OPEN QUESTION FOR \S4, NOT RESOLVED HERE.
%    (Roll)'s side condition reads "mu = A(phi U_I psi), sup I < infty;
%    otherwise mu^bullet is left unchanged".  What happens to the
%    re-demanded mu in the "otherwise" case is not stated: if it stays in
%    Gamma then (Cov) cannot fire on it (its side condition asks
%    O(mu) undefined) and the node never saturates.  \S5 is written so it
%    does not depend on the answer, using only that at most one pursuit
%    per modality is live.  The implementation discharges the re-demand
%    against the live pursuit (apply_cover, :117-144), which is
%    presumably the intent and should be said in \S4.
%
% 5. THE FIGURE.
%    sec/evaluation_new.tex cites \ref{fig:punctual} in "What this does
%    not show".  The figure is kept here and in the shipped section it
%    was dropped for space; either restore it in \S5 (costs ~0.25 page)
%    or move the float above into sec/evaluation_new.tex.
   % only the part below
%% =====================================================================
\iffalse
\section{Deferred proofs}\label{app:proofs}

\begin{proof}[of Lemma~\ref{lem:finiteness}]
$\Gamma$, $\Ctx$, the reset set and the discharge marker range over subsets of
$\FL^{\bullet}(\varphi)$ or of $\clockSet$; $\obligations$ is fixed by its domain; and $\Deadline$ gives each
$x_\mu$ one of the bounds $\sup\kI$, $\sup\kI-\inf\kI$ with a strictness flag. Each contributes
$2^{O(n)}$. A canonical $\cong_M$-zone over $K\le n+1$ clocks is a $(K{+}1)^{2}$ DBM with entries
$\infty$ or in $\{-M,\dots,M\}\times\{<,\le\}$, giving $(4M{+}3)^{(n+2)^{2}}=M^{O(n^{2})}$ zones.
Every rule has at most three children except \ruleref{EX}, with one per $\EX\psi\in\Gamma$.
\end{proof}

\begin{proof}[of Proposition~\ref{prop:gfp}]
\ref{D1} ignores $A$; a node kept by \ref{D2} has its witnessing child, or all its children, in $A$
and hence in any superset; \ref{D3} and \ref{D4} ask for a path in $A$, which survives enlargement.
So $F$ is monotone on the finite lattice $2^{G(\varphi)}$, and $A\subseteq F(A)$ unfolds to the
clauses defining a successful marking. Knaster--Tarski gives a greatest post-fixpoint, equal to $\nu
F$ and to the union of all post-fixpoints, and the tableau succeeds iff one of them contains $N_0$.
\end{proof}

\begin{proof}[of Lemma~\ref{lem:search}]
(i) The first claim is \ref{D2}. The explored subgraph is closed under $\Expand{}$ except at retired
nodes, so a path of $\mathcal M^{*}$ using an unexpanded child of $v$ passes through $v$; if
$v\notin\mathcal M^{*}$ there is no such path, and if $v\in\mathcal M^{*}$ then $v$ satisfies
\ref{M1}, and \ref{M2} for each pursuit active at it, so a node reaching $v$ inherits both witnesses
by prefixing the path to $v$, which carries a pursuit exactly when it is active at $v$.

(ii) $(\Leftarrow)$ Concatenating paths in $T$ gives a cycle through the \ruleref{Delay}-successor,
which therefore elapses positive time, and through one reset node of every clock $T$ resets; the rest
are unbounded throughout. $(\Rightarrow)$ A time-progressing cycle in $T$ elapses positive time,
hence crosses a \ruleref{Delay} step; and a clock reset nowhere in $T$ is forced unbounded on that
cycle, whence unbounded on all of $T$, since a path inside $T$ never resets it and its value, already
above $M$, only grows.

(iii) $C$ being maximal, a path of $\mathcal M^{*}$ from $N\in C$ either stays in $C$, and a witness
for \ref{M1} or \ref{M2} then lies inside $C$; or it leaves $C$ once and for all at some
$c\in\mathcal M^{*}\setminus C$, which by hypothesis lies in $\Out_C$ and itself satisfies \ref{M1},
and \ref{M2} for every pursuit active at it, so prefixing the path from $N$ to $c$ turns either
witness into one for $N$. Hence $\mathcal M^{*}\cap C\subseteq F(\mathcal M^{*}\cap C\cup\Out_C)$.
Conversely any $A\subseteq C$ with $A\subseteq F(A\cup\Out_C)$ extends by the nodes of
$\mathcal M^{*}$ reachable from $\Out_C$ to a post-fixpoint of $F$, which Proposition~\ref{prop:gfp}
places inside $\mathcal M^{*}$. The greatest elements coincide.
\end{proof}

\begin{proof}[of Lemma~\ref{lem:zones}]
(i) Every zone operation a rule performs (intersection with a single-clock guard, reset,
time-successor, $\cong_M$-canonicalisation) back-propagates, non-empty precisely because the
successor valuation lies in its image~\cite{Bouyer05TimedAutomata}; thread backwards from $\nu_k$,
noting that $\Dopen Z$ and $\Dbd Z$ lie strictly ahead of $Z$ in time. (ii) No clock crosses an
integer at or below $M$, so no integer part changes and no fractional ordering wraps, while clocks
above $M$ are identified by extrapolation; the segment lies in one region, on which TCTL truth is
constant~\cite{AlurCourcoubetisDill1990,Bouyer05TimedAutomata}, and its endpoints are nodes of their
own by the integer split of \ruleref{Delay}. (iii) A cycle elapsing no time takes zero time over
infinitely many traversals; otherwise some clock is neither reset nor extrapolated above $M$, so it
keeps a bound $c\le M$, grows by exactly the delay elapsed, and caps the total at $c$. Conversely (i)
yields a traversal elapsing some $\delta>0$, and the cycle returning to the same node, no rule
distinguishes two visits and no budget depletes, so $\delta$ recurs at every iteration.
\end{proof}

\begin{proof}[of Lemma~\ref{lem:recurrence}]
No rule drops $\mu^{\ast}$: \ruleref{Ustoproll} keeps it and \ruleref{Ugo} moves it to $\Ctx$, whence
\ruleref{EX} returns it to $\Gamma$ at every successor. It is carried forever, and with it the
deadline on $x_\mu$, which is $b$ until the first stop and $b-a$ after. Both are at most $M$, so
$x_\mu$ is bounded along $\rho$ and, by Lemma~\ref{lem:search}(ii), every time-progressing cycle
$\rho$ meets resets it. Only \ruleref{Ustoproll} does, \ruleref{Cov} being blocked while the pursuit
is live, and by \ref{M1} the schedule reaches such a cycle and traverses it without end, so the stop
recurs; its first guard is $x_\mu\in\kI$ and each later one is taken under $x_\mu\le b-a$ measured
from the preceding stop, which are the two bounds. Now let $t\ge t_0$ be a position and
$w_j$ the first stop at or after $t+a$. If $j=1$ then $w_1\le t_0+b\le t+b$; otherwise $w_{j-1}<t+a$,
so $w_j\le w_{j-1}+(b-a)<t+b$. Either way $w_j-t\in\kI$, and $\psi$ holds at $w_j$ while
\ruleref{Ugo} asserts $\varphi$ throughout, which is \eqref{eq:sat_u}.
\end{proof}
\fi


%% =====================================================================
%% DESIGN NOTES (not for the paper)
%% =====================================================================
%
% 1. WHY THE OLD SECTION COULD NOT BE EDITED IN PLACE.
%    sec/section5_OTF_old.tex is written against an earlier calculus: it
%    cites \ruleref{Covstar} and the satisfaction sets $\Sat_\Ctx(\mu)$.
%    Neither exists in sec/section4_nosat.tex, whose cover rule is the
%    conjunctive discharge/launch pair.  Lemma lem:sat of the old file,
%    its longest, related $\Sat_\Ctx$ to truth and has no referent any
%    more; its work is now done by the OR/AND case analysis inside
%    Theorem thm:completeness and by Lemma lem:recurrence.  The
%    $\Sat_\Ctx$ paragraph of the old complexity proof is likewise gone.
%
% 2. WHAT IS NEW RELATIVE TO THE OLD SECTION.
%    - Definition def:graph and the expansion order.  The old text never
%      said what the children of a node are; without a deterministic
%      \Expand{} the algorithm is not well defined and lem:finiteness
%      counts nothing.
%    - The observation that no local-termination lemma is needed:
%      identifying nodes up to \cong_M turns a regressing rule chain into
%      a cycle, and a cycle with no (Delay) step fails (M1).  This is why
%      the calculus can afford rules that re-impose their own premise
%      ((rec) of (U*), and (R^split)).
%    - Proposition prop:gfp, stated and proved rather than asserted.
%    - Lemma lem:recurrence.  This argument existed only as a
%      commented-out paragraph in section4_nosat.tex (lines 489-495).  It
%      is the semantic justification of the b-a retuning in (rec).
%      Without it, soundness for a persistent universal until is not
%      proved.
%    - The schedule / finite memory in the model extraction.  The old
%      proof said "its locations are the nodes of M*", which is wrong for
%      an OR-node with two surviving children: the automaton would
%      contain both, and a universal obligation could fail along the one
%      the marking did not use to establish the disjunction.  The
%      round-robin schedule is the standard fulfilling-path fix.
%    - Completeness restructured as R \subseteq F(R) for R the realisable
%      nodes.  Shorter, and it finally uses the realisability definition
%      of \S4.1, which the old proof did not.
%
% 3. FAITHFULNESS TO Implementation/src/tableau.cpp.
%    - Explore matches Engine::decide_scc (:1543); the settles test with
%      the top-of-stack guard is settles (:1571).
%    - \nu F_C matches evaluate_scc (:1284): optimistic start, inner
%      local-test fixpoint, then query_progress (:1392) and query_until
%      (:1435), outer loop until no deletion.
%    - Lemma lem:search(ii) is progressing() (:1368): union of reset
%      clocks, intersection of unbounded clocks, at least one (Delay).
%    - Back edges get a lowlink update and no verdict (:1584).  The notes
%      record this as load-bearing: optimistic back edges produced 82
%      soundness violations.
%    - "only the clocks of live pursuits and deadlines" is
%      collapse_dead_clocks (:698), an exact active-clock reduction.
%      (That sentence was cut from the shipped version for space.)
%    - The reset/discharge markers in node identity are just_reset and
%      discharges_until in acc_key (:1149-1157).
%    NOT claimed here, deliberately: the prototype still computes
%    satisfaction sets S_X for its cover rule and never populates
%    Node::X.  Both are artefacts of the older calculus.  Align the
%    prototype with \S4 before the evaluation numbers are cited.
%
% 4. OPEN QUESTION FOR \S4, NOT RESOLVED HERE.
%    (Roll)'s side condition reads "mu = A(phi U_I psi), sup I < infty;
%    otherwise mu^bullet is left unchanged".  What happens to the
%    re-demanded mu in the "otherwise" case is not stated: if it stays in
%    Gamma then (Cov) cannot fire on it (its side condition asks
%    O(mu) undefined) and the node never saturates.  \S5 is written so it
%    does not depend on the answer, using only that at most one pursuit
%    per modality is live.  The implementation discharges the re-demand
%    against the live pursuit (apply_cover, :117-144), which is
%    presumably the intent and should be said in \S4.
%
% 5. THE FIGURE.
%    sec/evaluation_new.tex cites \ref{fig:punctual} in "What this does
%    not show".  The figure is kept here and in the shipped section it
%    was dropped for space; either restore it in \S5 (costs ~0.25 page)
%    or move the float above into sec/evaluation_new.tex.
   % only the part below
%% =====================================================================
\iffalse
\section{Deferred proofs}\label{app:proofs}

\begin{proof}[of Lemma~\ref{lem:finiteness}]
$\Gamma$, $\Ctx$, the reset set and the discharge marker range over subsets of
$\FL^{\bullet}(\varphi)$ or of $\clockSet$; $\obligations$ is fixed by its domain; and $\Deadline$ gives each
$x_\mu$ one of the bounds $\sup\kI$, $\sup\kI-\inf\kI$ with a strictness flag. Each contributes
$2^{O(n)}$. A canonical $\cong_M$-zone over $K\le n+1$ clocks is a $(K{+}1)^{2}$ DBM with entries
$\infty$ or in $\{-M,\dots,M\}\times\{<,\le\}$, giving $(4M{+}3)^{(n+2)^{2}}=M^{O(n^{2})}$ zones.
Every rule has at most three children except \ruleref{EX}, with one per $\EX\psi\in\Gamma$.
\end{proof}

\begin{proof}[of Proposition~\ref{prop:gfp}]
\ref{D1} ignores $A$; a node kept by \ref{D2} has its witnessing child, or all its children, in $A$
and hence in any superset; \ref{D3} and \ref{D4} ask for a path in $A$, which survives enlargement.
So $F$ is monotone on the finite lattice $2^{G(\varphi)}$, and $A\subseteq F(A)$ unfolds to the
clauses defining a successful marking. Knaster--Tarski gives a greatest post-fixpoint, equal to $\nu
F$ and to the union of all post-fixpoints, and the tableau succeeds iff one of them contains $N_0$.
\end{proof}

\begin{proof}[of Lemma~\ref{lem:search}]
(i) The first claim is \ref{D2}. The explored subgraph is closed under $\Expand{}$ except at retired
nodes, so a path of $\mathcal M^{*}$ using an unexpanded child of $v$ passes through $v$; if
$v\notin\mathcal M^{*}$ there is no such path, and if $v\in\mathcal M^{*}$ then $v$ satisfies
\ref{M1}, and \ref{M2} for each pursuit active at it, so a node reaching $v$ inherits both witnesses
by prefixing the path to $v$, which carries a pursuit exactly when it is active at $v$.

(ii) $(\Leftarrow)$ Concatenating paths in $T$ gives a cycle through the \ruleref{Delay}-successor,
which therefore elapses positive time, and through one reset node of every clock $T$ resets; the rest
are unbounded throughout. $(\Rightarrow)$ A time-progressing cycle in $T$ elapses positive time,
hence crosses a \ruleref{Delay} step; and a clock reset nowhere in $T$ is forced unbounded on that
cycle, whence unbounded on all of $T$, since a path inside $T$ never resets it and its value, already
above $M$, only grows.

(iii) $C$ being maximal, a path of $\mathcal M^{*}$ from $N\in C$ either stays in $C$, and a witness
for \ref{M1} or \ref{M2} then lies inside $C$; or it leaves $C$ once and for all at some
$c\in\mathcal M^{*}\setminus C$, which by hypothesis lies in $\Out_C$ and itself satisfies \ref{M1},
and \ref{M2} for every pursuit active at it, so prefixing the path from $N$ to $c$ turns either
witness into one for $N$. Hence $\mathcal M^{*}\cap C\subseteq F(\mathcal M^{*}\cap C\cup\Out_C)$.
Conversely any $A\subseteq C$ with $A\subseteq F(A\cup\Out_C)$ extends by the nodes of
$\mathcal M^{*}$ reachable from $\Out_C$ to a post-fixpoint of $F$, which Proposition~\ref{prop:gfp}
places inside $\mathcal M^{*}$. The greatest elements coincide.
\end{proof}

\begin{proof}[of Lemma~\ref{lem:zones}]
(i) Every zone operation a rule performs (intersection with a single-clock guard, reset,
time-successor, $\cong_M$-canonicalisation) back-propagates, non-empty precisely because the
successor valuation lies in its image~\cite{Bouyer05TimedAutomata}; thread backwards from $\nu_k$,
noting that $\Dopen Z$ and $\Dbd Z$ lie strictly ahead of $Z$ in time. (ii) No clock crosses an
integer at or below $M$, so no integer part changes and no fractional ordering wraps, while clocks
above $M$ are identified by extrapolation; the segment lies in one region, on which TCTL truth is
constant~\cite{AlurCourcoubetisDill1990,Bouyer05TimedAutomata}, and its endpoints are nodes of their
own by the integer split of \ruleref{Delay}. (iii) A cycle elapsing no time takes zero time over
infinitely many traversals; otherwise some clock is neither reset nor extrapolated above $M$, so it
keeps a bound $c\le M$, grows by exactly the delay elapsed, and caps the total at $c$. Conversely (i)
yields a traversal elapsing some $\delta>0$, and the cycle returning to the same node, no rule
distinguishes two visits and no budget depletes, so $\delta$ recurs at every iteration.
\end{proof}

\begin{proof}[of Lemma~\ref{lem:recurrence}]
No rule drops $\mu^{\ast}$: \ruleref{Ustoproll} keeps it and \ruleref{Ugo} moves it to $\Ctx$, whence
\ruleref{EX} returns it to $\Gamma$ at every successor. It is carried forever, and with it the
deadline on $x_\mu$, which is $b$ until the first stop and $b-a$ after. Both are at most $M$, so
$x_\mu$ is bounded along $\rho$ and, by Lemma~\ref{lem:search}(ii), every time-progressing cycle
$\rho$ meets resets it. Only \ruleref{Ustoproll} does, \ruleref{Cov} being blocked while the pursuit
is live, and by \ref{M1} the schedule reaches such a cycle and traverses it without end, so the stop
recurs; its first guard is $x_\mu\in\kI$ and each later one is taken under $x_\mu\le b-a$ measured
from the preceding stop, which are the two bounds. Now let $t\ge t_0$ be a position and
$w_j$ the first stop at or after $t+a$. If $j=1$ then $w_1\le t_0+b\le t+b$; otherwise $w_{j-1}<t+a$,
so $w_j\le w_{j-1}+(b-a)<t+b$. Either way $w_j-t\in\kI$, and $\psi$ holds at $w_j$ while
\ruleref{Ugo} asserts $\varphi$ throughout, which is \eqref{eq:sat_u}.
\end{proof}
\fi


%% =====================================================================
%% DESIGN NOTES (not for the paper)
%% =====================================================================
%
% 1. WHY THE OLD SECTION COULD NOT BE EDITED IN PLACE.
%    sec/section5_OTF_old.tex is written against an earlier calculus: it
%    cites \ruleref{Covstar} and the satisfaction sets $\Sat_\Ctx(\mu)$.
%    Neither exists in sec/section4_nosat.tex, whose cover rule is the
%    conjunctive discharge/launch pair.  Lemma lem:sat of the old file,
%    its longest, related $\Sat_\Ctx$ to truth and has no referent any
%    more; its work is now done by the OR/AND case analysis inside
%    Theorem thm:completeness and by Lemma lem:recurrence.  The
%    $\Sat_\Ctx$ paragraph of the old complexity proof is likewise gone.
%
% 2. WHAT IS NEW RELATIVE TO THE OLD SECTION.
%    - Definition def:graph and the expansion order.  The old text never
%      said what the children of a node are; without a deterministic
%      \Expand{} the algorithm is not well defined and lem:finiteness
%      counts nothing.
%    - The observation that no local-termination lemma is needed:
%      identifying nodes up to \cong_M turns a regressing rule chain into
%      a cycle, and a cycle with no (Delay) step fails (M1).  This is why
%      the calculus can afford rules that re-impose their own premise
%      ((rec) of (U*), and (R^split)).
%    - Proposition prop:gfp, stated and proved rather than asserted.
%    - Lemma lem:recurrence.  This argument existed only as a
%      commented-out paragraph in section4_nosat.tex (lines 489-495).  It
%      is the semantic justification of the b-a retuning in (rec).
%      Without it, soundness for a persistent universal until is not
%      proved.
%    - The schedule / finite memory in the model extraction.  The old
%      proof said "its locations are the nodes of M*", which is wrong for
%      an OR-node with two surviving children: the automaton would
%      contain both, and a universal obligation could fail along the one
%      the marking did not use to establish the disjunction.  The
%      round-robin schedule is the standard fulfilling-path fix.
%    - Completeness restructured as R \subseteq F(R) for R the realisable
%      nodes.  Shorter, and it finally uses the realisability definition
%      of \S4.1, which the old proof did not.
%
% 3. FAITHFULNESS TO Implementation/src/tableau.cpp.
%    - Explore matches Engine::decide_scc (:1543); the settles test with
%      the top-of-stack guard is settles (:1571).
%    - \nu F_C matches evaluate_scc (:1284): optimistic start, inner
%      local-test fixpoint, then query_progress (:1392) and query_until
%      (:1435), outer loop until no deletion.
%    - Lemma lem:search(ii) is progressing() (:1368): union of reset
%      clocks, intersection of unbounded clocks, at least one (Delay).
%    - Back edges get a lowlink update and no verdict (:1584).  The notes
%      record this as load-bearing: optimistic back edges produced 82
%      soundness violations.
%    - "only the clocks of live pursuits and deadlines" is
%      collapse_dead_clocks (:698), an exact active-clock reduction.
%      (That sentence was cut from the shipped version for space.)
%    - The reset/discharge markers in node identity are just_reset and
%      discharges_until in acc_key (:1149-1157).
%    NOT claimed here, deliberately: the prototype still computes
%    satisfaction sets S_X for its cover rule and never populates
%    Node::X.  Both are artefacts of the older calculus.  Align the
%    prototype with \S4 before the evaluation numbers are cited.
%
% 4. OPEN QUESTION FOR \S4, NOT RESOLVED HERE.
%    (Roll)'s side condition reads "mu = A(phi U_I psi), sup I < infty;
%    otherwise mu^bullet is left unchanged".  What happens to the
%    re-demanded mu in the "otherwise" case is not stated: if it stays in
%    Gamma then (Cov) cannot fire on it (its side condition asks
%    O(mu) undefined) and the node never saturates.  \S5 is written so it
%    does not depend on the answer, using only that at most one pursuit
%    per modality is live.  The implementation discharges the re-demand
%    against the live pursuit (apply_cover, :117-144), which is
%    presumably the intent and should be said in \S4.
%
% 5. THE FIGURE.
%    sec/evaluation_new.tex cites \ref{fig:punctual} in "What this does
%    not show".  The figure is kept here and in the shipped section it
%    was dropped for space; either restore it in \S5 (costs ~0.25 page)
%    or move the float above into sec/evaluation_new.tex.
   % only the part below
%% =====================================================================
\iffalse
\section{Deferred proofs}\label{app:proofs}

\begin{proof}[of Lemma~\ref{lem:finiteness}]
$\Gamma$, $\Ctx$, the reset set and the discharge marker range over subsets of
$\FL^{\bullet}(\varphi)$ or of $\clockSet$; $\obligations$ is fixed by its domain; and $\Deadline$ gives each
$x_\mu$ one of the bounds $\sup\kI$, $\sup\kI-\inf\kI$ with a strictness flag. Each contributes
$2^{O(n)}$. A canonical $\cong_M$-zone over $K\le n+1$ clocks is a $(K{+}1)^{2}$ DBM with entries
$\infty$ or in $\{-M,\dots,M\}\times\{<,\le\}$, giving $(4M{+}3)^{(n+2)^{2}}=M^{O(n^{2})}$ zones.
Every rule has at most three children except \ruleref{EX}, with one per $\EX\psi\in\Gamma$.
\end{proof}

\begin{proof}[of Proposition~\ref{prop:gfp}]
\ref{D1} ignores $A$; a node kept by \ref{D2} has its witnessing child, or all its children, in $A$
and hence in any superset; \ref{D3} and \ref{D4} ask for a path in $A$, which survives enlargement.
So $F$ is monotone on the finite lattice $2^{G(\varphi)}$, and $A\subseteq F(A)$ unfolds to the
clauses defining a successful marking. Knaster--Tarski gives a greatest post-fixpoint, equal to $\nu
F$ and to the union of all post-fixpoints, and the tableau succeeds iff one of them contains $N_0$.
\end{proof}

\begin{proof}[of Lemma~\ref{lem:search}]
(i) The first claim is \ref{D2}. The explored subgraph is closed under $\Expand{}$ except at retired
nodes, so a path of $\mathcal M^{*}$ using an unexpanded child of $v$ passes through $v$; if
$v\notin\mathcal M^{*}$ there is no such path, and if $v\in\mathcal M^{*}$ then $v$ satisfies
\ref{M1}, and \ref{M2} for each pursuit active at it, so a node reaching $v$ inherits both witnesses
by prefixing the path to $v$, which carries a pursuit exactly when it is active at $v$.

(ii) $(\Leftarrow)$ Concatenating paths in $T$ gives a cycle through the \ruleref{Delay}-successor,
which therefore elapses positive time, and through one reset node of every clock $T$ resets; the rest
are unbounded throughout. $(\Rightarrow)$ A time-progressing cycle in $T$ elapses positive time,
hence crosses a \ruleref{Delay} step; and a clock reset nowhere in $T$ is forced unbounded on that
cycle, whence unbounded on all of $T$, since a path inside $T$ never resets it and its value, already
above $M$, only grows.

(iii) $C$ being maximal, a path of $\mathcal M^{*}$ from $N\in C$ either stays in $C$, and a witness
for \ref{M1} or \ref{M2} then lies inside $C$; or it leaves $C$ once and for all at some
$c\in\mathcal M^{*}\setminus C$, which by hypothesis lies in $\Out_C$ and itself satisfies \ref{M1},
and \ref{M2} for every pursuit active at it, so prefixing the path from $N$ to $c$ turns either
witness into one for $N$. Hence $\mathcal M^{*}\cap C\subseteq F(\mathcal M^{*}\cap C\cup\Out_C)$.
Conversely any $A\subseteq C$ with $A\subseteq F(A\cup\Out_C)$ extends by the nodes of
$\mathcal M^{*}$ reachable from $\Out_C$ to a post-fixpoint of $F$, which Proposition~\ref{prop:gfp}
places inside $\mathcal M^{*}$. The greatest elements coincide.
\end{proof}

\begin{proof}[of Lemma~\ref{lem:zones}]
(i) Every zone operation a rule performs (intersection with a single-clock guard, reset,
time-successor, $\cong_M$-canonicalisation) back-propagates, non-empty precisely because the
successor valuation lies in its image~\cite{Bouyer05TimedAutomata}; thread backwards from $\nu_k$,
noting that $\Dopen Z$ and $\Dbd Z$ lie strictly ahead of $Z$ in time. (ii) No clock crosses an
integer at or below $M$, so no integer part changes and no fractional ordering wraps, while clocks
above $M$ are identified by extrapolation; the segment lies in one region, on which TCTL truth is
constant~\cite{AlurCourcoubetisDill1990,Bouyer05TimedAutomata}, and its endpoints are nodes of their
own by the integer split of \ruleref{Delay}. (iii) A cycle elapsing no time takes zero time over
infinitely many traversals; otherwise some clock is neither reset nor extrapolated above $M$, so it
keeps a bound $c\le M$, grows by exactly the delay elapsed, and caps the total at $c$. Conversely (i)
yields a traversal elapsing some $\delta>0$, and the cycle returning to the same node, no rule
distinguishes two visits and no budget depletes, so $\delta$ recurs at every iteration.
\end{proof}

\begin{proof}[of Lemma~\ref{lem:recurrence}]
No rule drops $\mu^{\ast}$: \ruleref{Ustoproll} keeps it and \ruleref{Ugo} moves it to $\Ctx$, whence
\ruleref{EX} returns it to $\Gamma$ at every successor. It is carried forever, and with it the
deadline on $x_\mu$, which is $b$ until the first stop and $b-a$ after. Both are at most $M$, so
$x_\mu$ is bounded along $\rho$ and, by Lemma~\ref{lem:search}(ii), every time-progressing cycle
$\rho$ meets resets it. Only \ruleref{Ustoproll} does, \ruleref{Cov} being blocked while the pursuit
is live, and by \ref{M1} the schedule reaches such a cycle and traverses it without end, so the stop
recurs; its first guard is $x_\mu\in\kI$ and each later one is taken under $x_\mu\le b-a$ measured
from the preceding stop, which are the two bounds. Now let $t\ge t_0$ be a position and
$w_j$ the first stop at or after $t+a$. If $j=1$ then $w_1\le t_0+b\le t+b$; otherwise $w_{j-1}<t+a$,
so $w_j\le w_{j-1}+(b-a)<t+b$. Either way $w_j-t\in\kI$, and $\psi$ holds at $w_j$ while
\ruleref{Ugo} asserts $\varphi$ throughout, which is \eqref{eq:sat_u}.
\end{proof}
\fi


%% =====================================================================
%% DESIGN NOTES (not for the paper)
%% =====================================================================
%
% 1. WHY THE OLD SECTION COULD NOT BE EDITED IN PLACE.
%    sec/section5_OTF_old.tex is written against an earlier calculus: it
%    cites \ruleref{Covstar} and the satisfaction sets $\Sat_\Ctx(\mu)$.
%    Neither exists in sec/section4_nosat.tex, whose cover rule is the
%    conjunctive discharge/launch pair.  Lemma lem:sat of the old file,
%    its longest, related $\Sat_\Ctx$ to truth and has no referent any
%    more; its work is now done by the OR/AND case analysis inside
%    Theorem thm:completeness and by Lemma lem:recurrence.  The
%    $\Sat_\Ctx$ paragraph of the old complexity proof is likewise gone.
%
% 2. WHAT IS NEW RELATIVE TO THE OLD SECTION.
%    - Definition def:graph and the expansion order.  The old text never
%      said what the children of a node are; without a deterministic
%      \Expand{} the algorithm is not well defined and lem:finiteness
%      counts nothing.
%    - The observation that no local-termination lemma is needed:
%      identifying nodes up to \cong_M turns a regressing rule chain into
%      a cycle, and a cycle with no (Delay) step fails (M1).  This is why
%      the calculus can afford rules that re-impose their own premise
%      ((rec) of (U*), and (R^split)).
%    - Proposition prop:gfp, stated and proved rather than asserted.
%    - Lemma lem:recurrence.  This argument existed only as a
%      commented-out paragraph in section4_nosat.tex (lines 489-495).  It
%      is the semantic justification of the b-a retuning in (rec).
%      Without it, soundness for a persistent universal until is not
%      proved.
%    - The schedule / finite memory in the model extraction.  The old
%      proof said "its locations are the nodes of M*", which is wrong for
%      an OR-node with two surviving children: the automaton would
%      contain both, and a universal obligation could fail along the one
%      the marking did not use to establish the disjunction.  The
%      round-robin schedule is the standard fulfilling-path fix.
%    - Completeness restructured as R \subseteq F(R) for R the realisable
%      nodes.  Shorter, and it finally uses the realisability definition
%      of \S4.1, which the old proof did not.
%
% 3. FAITHFULNESS TO Implementation/src/tableau.cpp.
%    - Explore matches Engine::decide_scc (:1543); the settles test with
%      the top-of-stack guard is settles (:1571).
%    - \nu F_C matches evaluate_scc (:1284): optimistic start, inner
%      local-test fixpoint, then query_progress (:1392) and query_until
%      (:1435), outer loop until no deletion.
%    - Lemma lem:search(ii) is progressing() (:1368): union of reset
%      clocks, intersection of unbounded clocks, at least one (Delay).
%    - Back edges get a lowlink update and no verdict (:1584).  The notes
%      record this as load-bearing: optimistic back edges produced 82
%      soundness violations.
%    - "only the clocks of live pursuits and deadlines" is
%      collapse_dead_clocks (:698), an exact active-clock reduction.
%      (That sentence was cut from the shipped version for space.)
%    - The reset/discharge markers in node identity are just_reset and
%      discharges_until in acc_key (:1149-1157).
%    NOT claimed here, deliberately: the prototype still computes
%    satisfaction sets S_X for its cover rule and never populates
%    Node::X.  Both are artefacts of the older calculus.  Align the
%    prototype with \S4 before the evaluation numbers are cited.
%
% 4. OPEN QUESTION FOR \S4, NOT RESOLVED HERE.
%    (Roll)'s side condition reads "mu = A(phi U_I psi), sup I < infty;
%    otherwise mu^bullet is left unchanged".  What happens to the
%    re-demanded mu in the "otherwise" case is not stated: if it stays in
%    Gamma then (Cov) cannot fire on it (its side condition asks
%    O(mu) undefined) and the node never saturates.  \S5 is written so it
%    does not depend on the answer, using only that at most one pursuit
%    per modality is live.  The implementation discharges the re-demand
%    against the live pursuit (apply_cover, :117-144), which is
%    presumably the intent and should be said in \S4.
%
% 5. THE FIGURE.
%    sec/evaluation_new.tex cites \ref{fig:punctual} in "What this does
%    not show".  The figure is kept here and in the shipped section it
%    was dropped for space; either restore it in \S5 (costs ~0.25 page)
%    or move the float above into sec/evaluation_new.tex.
